Each term is defined where it first appears; the chapter is given in brackets.
\begin{description}[style=nextline,leftmargin=1.2em,font=\normalfont\bfseries]
\item[1 by 1 convolution] A convolution that looks at one position at a time and mixes only channels: a linear layer applied at every pixel with shared weights. (\cref{ch:decoder})
\item[1-pixel tolerance] The rule that a flagged pixel adjacent to an object also catches it, allowing for ground truth about one pixel off the grid. (\cref{ch:measures})
\item[1x1 convolution] A convolution with a one-pixel kernel: a linear layer applied independently at every pixel, mixing channels only. (\cref{ch:shapes})
\item[3D ROC] Plots of detection and false-alarm rates against a min-max normalised threshold; its areas reduce to mean normalised scores and are degenerate for heavy-tailed scores. (\cref{ch:measures})
\item[4-connectivity] Pixels are neighbours only if they share an edge (up, down, left, right), not a corner. (\cref{ch:measures})
\item[40 percent rule] The trainer's rule that drops any patch with less than 40 percent valid pixels before the loss is computed. (\cref{ch:losses})
\item[ABU] The Airport-Beach-Urban benchmark of 13 small scenes (2017); radiance of unknown gain and no published wavelengths. (\cref{ch:benchmarks})
\item[Abundance] The fraction of a pixel's footprint covered by a material; the a of the linear mixing model. (\cref{ch:benchmarks})
\item[Accuracy] The fraction of answers a yes/no rule gets right; meaningless when one class is extremely rare. (\cref{ch:problem})
\item[Action] One processing step inside a Project; reads earlier Actions' outputs and writes its own to its own folder. (\cref{ch:product})
\item[Action chain] The sequence of product actions from upload to export: band filter, segmentation, scoring, detection prep, spectral match, export. (\cref{ch:inference})
\item[ActionOutput] The folder of files one Action writes, referred to by later Actions. (\cref{ch:product})
\item[Activation function] A nonlinear function applied after a linear layer, such as GELU or ReLU. (\cref{ch:shapes})
\item[Activation memory] Memory held by intermediate tensors of the forward pass that the backward pass needs; for the per-band stem about 19 tensors of 202 MiB per patch in bf16. (\cref{ch:crossstem})
\item[AdaBN] Adaptive batch normalisation: replacing BatchNorm's stored statistics with statistics measured on the data being scored. (\cref{ch:finetune})
\item[Adam] The optimiser that divides an EMA of each parameter's gradient by the square root of an EMA of its squared gradient, giving every parameter a step of about the learning rate when its gradient is consistent. (\cref{ch:trainloop})
\item[AdamW] Adam with decoupled weight decay: the decay is subtracted separately rather than passed through Adam's per-parameter normaliser. (\cref{ch:trainloop})
\item[Adapt stage] Fine-tuning reconstruction only, on all three sensors at their own bands, at a low learning rate. (\cref{ch:finetune})
\item[Adjusted Rand index] A pair-counting agreement between two groupings, corrected so a random grouping scores 0 and an identical one 1. (\cref{ch:hslwhat})
\item[Affine map (LayerNorm)] The learned per-feature scale gamma and shift beta applied after standardisation. (\cref{ch:block})
\item[Affine subspace] A subspace shifted by a fixed vector; Indradhanu's reconstructed spectra lie in a 32-dimensional one. (\cref{ch:compressor})
\item[Affine transform] Six numbers mapping a pixel's column and row to map coordinates: X = a col + b row + c, Y = d col + e row + f. (\cref{ch:material})
\item[Aliasing] Two different inputs giving the same encoding because a wave has completed a whole number of periods between them. (\cref{ch:ceiling})
\item[align\_corners=False] PyTorch's resizing convention in which pixel centres sit at half-integer positions and the outer edges of input and output grids line up. (\cref{ch:decoder})
\item[All-MLP decoder] SegFormer's decoder design: per-scale 1 by 1 projections, upsampling to one size, concatenation and a 1 by 1 fuse; 1 by 1 convolutions are per-pixel MLP layers. (\cref{ch:decoder})
\item[Allotrope] The web application this book describes: users upload satellite scenes, run chains of analysis steps, pick a threshold and export a bundle of candidate anomalies. (\cref{ch:problem})
\item[Angle between vectors] The angle whose cosine is the dot product divided by both lengths; it measures shape and ignores brightness. (\cref{ch:linalg})
\item[Annulus] Another name for the dual-window background ring. (\cref{ch:localrx})
\item[Anomaly] A pixel or small group of pixels unlike the rest of its scene; the word says nothing about what the pixel is. (\cref{ch:problem})
\item[Anomaly mask] The binary raster (anomaly\_mask.tif) written when a threshold is committed: 1 where the composite is at or above the threshold. (\cref{ch:material})
\item[Anomaly score] A non-negative number per pixel, larger meaning the pixel fits the model of normal worse. (\cref{ch:problem})
\item[Antardhana] The convolutional autoencoder trained to fill 13 to 25 percent hidden pixels with an MSE loss; 264,961 parameters. (\cref{ch:convae})
\item[api] The web-server process: handles requests and the database, hands out work, never computes. (\cref{ch:product})
\item[Argsort] The permutation that sorts a list: its k-th entry is the index of the k-th smallest value. (\cref{ch:masking})
\item[Array] A block of numbers indexed by whole-number positions; its shape lists the size of each index. (\cref{ch:light})
\item[Asanskrita] The mask-aware convolutional autoencoder with no normalisation: three input channels (temperature, validity, input mask); 267,009 parameters. (\cref{ch:convae})
\item[atan2] The two-argument arctangent, giving the angle of the point (x, y) from the positive x-axis; used by SAMLoss to compute the spectral angle with finite derivatives. (\cref{ch:losses})
\item[Atmospheric window (absorption)] A wavelength range where water vapour or carbon dioxide absorbs so strongly that a band measures noise; removed from the grid. (\cref{ch:sensors})
\item[Attention mask] A boolean tensor passed to PyTorch's fused attention function; True means attention is allowed, False means the score becomes minus infinity. (\cref{ch:crossstem})
\item[Attention matrix] The N-by-N (or N-by-N-reduced) matrix of softmax weights; row i says how query i shares its attention among the keys, and every row sums to one. (\cref{ch:attention})
\item[Attention weights] The softmax of scaled query-key scores; each row is positive and sums to 1. (\cref{ch:ceiling})
\item[AUC] The area under the ROC curve; 1 is perfect, 0.5 is random. Equal to the probability that a random target outscores a random background pixel, ties counting half. (\cref{ch:measures})
\item[Autocast] PyTorch's context that runs eligible operations in a lower-precision format automatically. (\cref{ch:trainloop})
\item[Autoencoder] A reconstruction model built from an encoder that summarises the input and a decoder that rebuilds the input from the summary. (\cref{ch:recon})
\item[Average pooling] Replacing each block of a grid by the plain average of its entries; a special case of the learned reduction convolution. (\cref{ch:esa})
\item[AVIRIS-NG] Airborne hyperspectral sensor with about 425 bands, delivered as ENVI files in a folder. (\cref{ch:sensors})
\item[Axis] One direction along which a tensor is indexed; also called a dimension. (\cref{ch:shapes})
\item[Background subspace] The span of a scene's top ten principal directions. (\cref{ch:benchmarks})
\item[Backpropagation] Computing the gradient of the loss with respect to every parameter by applying the chain rule backwards from the loss through each layer. (\cref{ch:losses})
\item[Band] One measurement channel of a sensor, recording light in one narrow range of wavelengths; called a channel in deep-learning code. (\cref{ch:light})
\item[Band alignment] Matching a sample's bands to a target list by wavelength within 0.5 nm; unmatched targets become nodata and unmatched sources are dropped. (\cref{ch:pretrain})
\item[Band density] How many bands a sensor places in a region of the spectrum; with a one-vote-per-band average it changes the stem's output for the same spectrum. (\cref{ch:pretrain})
\item[Band grouping] Making one token from several adjacent bands to cut the token count (202 to 51), at the cost of mixing fine spectral detail early. (\cref{ch:crossstem})
\item[Band pipeline] The eight stages that turn a sensor's native bands into a complete spectrum per valid pixel on the 165-band common grid. (\cref{ch:bandpipeline})
\item[Band response] The weighting over wavelengths with which a sensor band integrates light, roughly Gaussian with a given full width at half maximum. (\cref{ch:material})
\item[Band token] One band of one pixel represented as a d-number vector: its value stretched by Linear(1, d), plus the encoding of its wavelength. (\cref{ch:ceiling})
\item[Band validity mask] The tensor \texttt{band\_valid}: one boolean per band per pixel, False for a dead or missing band. (\cref{ch:crossstem})
\item[Band weighting by width] An option of the reduced stem in which each band counts by the nanometres of spectrum it covers, making its means approximate integrals over wavelength. (\cref{ch:specdecoder})
\item[Band-delta compression] Storing each band as its int16 difference from the previous band, with wrap-around, then compressing with zstd; exact and about 2 times smaller. (\cref{ch:pretrain})
\item[Band-grid check] The worker's test that a model's expected grid matches the vendable's wavelengths, failing the action with a message naming the band-filter setting. (\cref{ch:material})
\item[Band-list augmentation] Training on a random subset of a batch's bands while reconstructing all of them; safe because the loss is taken only on hidden pixels. (\cref{ch:specdecoder})
\item[BandFilterConfig] The configuration object that drives the eight-stage band pipeline: exclusion windows, edge trimming, coverage threshold, voxel-fraction threshold, quality masks and target grid. (\cref{ch:bandpipeline})
\item[Base rate] The fraction of all cases that belong to the rare class, here the fraction of pixels that are anomalous. (\cref{ch:problem})
\item[Base width] The channel count of a convolutional autoencoder's first stage; later stages double it. (\cref{ch:convae})
\item[Batch] A group of image patches processed together in one step of training or inference; its size is written B. (\cref{ch:light})
\item[Batch (mini-batch)] The group of patches processed together in one forward pass; its gradient is a noisy estimate of the whole training set's gradient. (\cref{ch:trainloop})
\item[Batch normalisation] A layer that standardises each channel with the mean and variance computed over the batch and all pixels, then applies a learned scale and shift. (\cref{ch:compressor})
\item[BatchNorm] Normalisation of one channel using statistics gathered across the batch, with running statistics at evaluation; its output depends on the other samples in the batch. (\cref{ch:block})
\item[BatchNorm running statistics] A running mean and variance per channel and one batch counter, kept as buffers for evaluation. (\cref{ch:params})
\item[Bench] The mean ROC-AUC on the four labelled PRISMA scenes under residual Mahalanobis, the detection measure every fine-tuning gate reads. (\cref{ch:finetune})
\item[Bench resolution] The smallest difference a bench can tell from noise: about plus or minus 0.005 to 0.01 pAUC on the real PRISMA bench. (\cref{ch:bootstrap})
\item[bench-v1, bench-v2, bench-v2.1] The implant bench before foreignness and halo were matched; after (accepted 2026-10-17); and rebuilt after the descending-wavelength fix with a tighter foreignness tolerance (2026-10-18). (\cref{ch:benchmarks})
\item[bench-v2.1] The implant benchmark rebuilt on 2026-10-18 with sorted interpolation and foreignness tolerance 1.2: 18 hosts, 2,817 objects. (\cref{ch:sortbug})
\item[Bernoulli distribution] The distribution of a single yes/no outcome with probability p of yes: P(t) = $p^t(1-p)^{1-t}$. (\cref{ch:seghead})
\item[Bernoulli variable] A random variable that is 1 with probability p and 0 otherwise; its mean is p and its variance p(1 - p). (\cref{ch:sortbug})
\item[best.pt] The checkpoint a run keeps as its best, chosen by validation reconstruction (or segmentation) loss, not by detection. (\cref{ch:ftruns})
\item[bfloat16] A 16-bit floating-point format with float32's 8 exponent bits and 7 mantissa bits: float32's range at about 2 to 3 significant digits. (\cref{ch:trainloop})
\item[Bias] A learned constant added to a neuron's or channel's output. (\cref{ch:shapes})
\item[Bias correction] Dividing Adam's moment estimates by 1 minus beta to the power t, which removes the bias towards zero caused by starting the averages at zero. (\cref{ch:trainloop})
\item[Biased variance] A variance that divides the sum of squared deviations by the number of values M; BatchNorm normalises with it. (\cref{ch:compressor})
\item[BIL] Band-interleaved-by-line layout (H, C, W): PRISMA's native layout. (\cref{ch:sensors})
\item[Bilinear interpolation] Linear interpolation along the rows and then between the rows; each output is a weighted average of four neighbouring samples with weights that sum to 1. (\cref{ch:decoder})
\item[Bin token] The mean over a bin's valid bands of each band value times its own wavelength-computed vector k(lambda), plus a bias from the bin centre's tag. (\cref{ch:specdecoder})
\item[Binary cross-entropy] The negative log-likelihood of 0/1 labels under predicted probabilities, averaged over pixels; its gradient with respect to the logit is p - t. (\cref{ch:seghead})
\item[Binary cross-entropy (BCE)] The loss -(y ln p + (1 - y) ln(1 - p)) for a yes/no target y and predicted probability p. (\cref{ch:ftruns})
\item[Binary dilation] Adding every pixel within the radius of the set; grows the invalid region. Computed as a stride-1 max pool. (\cref{ch:masking})
\item[Binary erosion] Keeping a pixel only when its whole square neighbourhood is in the set; shrinks the valid region by the radius. (\cref{ch:masking})
\item[Binary search] Finding a value's position in a sorted list by repeated halving; requires ascending order, which is why wavelengths must be sorted before np.interp or np.searchsorted. (\cref{ch:bandpipeline})
\item[Binary validity] The state after stage 7 in which every pixel is valid in all bands or in none. (\cref{ch:bandpipeline})
\item[Binomial standard error] The spread of a share caught over n independent objects, the square root of p(1 - p)/n; a rough ruler for whether two shares differ. (\cref{ch:sortbug})
\item[BIP] Band-interleaved-by-pixel layout (H, W, C): each pixel's spectrum is contiguous. (\cref{ch:sensors})
\item[Bit flag] A single binary digit of an integer used as a yes/no answer; bit k of q is (q >> k) and 1. (\cref{ch:sensors})
\item[Blind spot] The arrangement at inference in which each pixel is reconstructed while its own token is hidden, so it cannot copy itself. (\cref{ch:ftruns})
\item[Blind unmixing] Finding both the endmembers and the abundances from the image alone. (\cref{ch:hslunmix})
\item[Block] In F001, 200 consecutive shards (30,000 patches, one epoch's worth) trained for 20 epochs before the next block replaces it. (\cref{ch:pretrain})
\item[Block masking] Hiding tokens in b x b groups rather than singly (FT08: 2 x 2 tokens, 8 px), so a target and its halo vanish together. (\cref{ch:decoderwork})
\item[Block parameter formula] One transformer block of width C, reduction ratio R and expansion 4 has 12 C squared + 53 C parameters, plus R squared C squared + 3C when R is above 1. (\cref{ch:params})
\item[Block-diagonal matrix] A square matrix whose non-zero entries sit in square blocks along the diagonal; its inverse is the inverse of each block in place. (\cref{ch:residualmaha})
\item[Boolean indexing] Selecting the entries of a tensor where a True/False tensor is True, producing a flat list; used to pull out the graded pixels. (\cref{ch:losses})
\item[Bootstrap] Estimating a statistic's sampling distribution by recomputing it on many samples drawn with replacement from the data. (\cref{ch:bootstrap})
\item[Bootstrap sample] One draw of n items with replacement from the n data items. (\cref{ch:bootstrap})
\item[Bottleneck] The internal summary between encoder and decoder; it constrains what the model can carry from input to output. (\cref{ch:recon})
\item[Box sum] The sum of a rectangle of values, obtained from an integral image with four look-ups: far corner minus top strip minus left strip plus the doubly removed corner. (\cref{ch:residualmaha})
\item[Broadcasting] Repeating an array along axes of size 1 so that arrays of shapes such as (B, 165, H, W) and (1, 165, 1, 1) combine element by element. (\cref{ch:bandpipeline})
\item[BSQ] Band-sequential layout (C, H, W): each band image is contiguous; the layout every model expects. (\cref{ch:sensors})
\item[Buffer] A tensor stored in a model's state but not trained, such as normalisation statistics or BatchNorm running averages; it travels inside the checkpoint. (\cref{ch:recon})
\item[Buffer (PyTorch)] A tensor registered on a module that is saved in the checkpoint and moved with the model but never trained; how the normalisation statistics are stored. (\cref{ch:bandpipeline})
\item[Cache key] The path of a cached map: dataset, scene, model id, and a SHA-1 of the scorer's settings joined to the scene fingerprint. (\cref{ch:runeval})
\item[Capacity] How much a model can represent; for fixed depth it grows with width. (\cref{ch:heavy})
\item[Capacity block] A GPU instance rented for a fixed, pre-booked time window that runs whether or not a job is running. (\cref{ch:heavy})
\item[Capacity limit] Error that remains even when the network can see the pixel, because it cannot render the detail. (\cref{ch:decoderwork})
\item[Catastrophic cancellation] Loss of precision when subtracting two large, nearly equal floating-point numbers, as in E[x squared] minus (E[x]) squared. (\cref{ch:bandpipeline})
\item[Centre wavelength] The wavelength at which a band responds most strongly; stored per band as band\_cw\_order. (\cref{ch:light})
\item[Centring] Subtracting the mean from every vector, so that the deviations add up to zero. (\cref{ch:linalg})
\item[Chain rule] The derivative of a composition is the product of the derivatives along the way: d g(f(x))/dx = g'(f(x)) f'(x). (\cref{ch:block})
\item[Chakshu] The thermal foundation model (``the eye''), the one-band sibling of Indradhanu. (\cref{ch:problem})
\item[Chance level] What a random scorer would get on a measure: 0.5 for AUC, half the limit for pAUC, the FPR itself for TPR at FPR, and 1 - (1 - q)\^{}m per object for objects caught. (\cref{ch:measures})
\item[Chance row] The first row of every report table: the mean over objects of each object's chance of being caught by random scores. (\cref{ch:runeval})
\item[Channel identity] A learned 32-number vector per latent channel, added to its stretched value, because channels have no wavelength to encode. (\cref{ch:specdecoder})
\item[Characteristic polynomial] det(S - lambda I) as a polynomial in lambda; its roots are the eigenvalues. (\cref{ch:linalg})
\item[Chebyshev distance] The larger of the row and column offsets, so a diagonal step counts as one. (\cref{ch:namer})
\item[Checkerboard cell size] The size, in stage-1 tokens, of the cells of the two complementary inference masks. (\cref{ch:decoderwork})
\item[Checkerboard mask] The default inference masking pattern, (row + column) mod 2 over the token grid; pass 1 hides the even tokens and pass 2 the odd ones. (\cref{ch:inference})
\item[Checkpoint] A saved file holding the model's parameters, the optimiser and scheduler state, the epoch and the configuration, from which training can resume. (\cref{ch:trainloop})
\item[Checkpoint noise] Variation in a measured score from one epoch's weights to the next, invisible to an object bootstrap. (\cref{ch:decoderwork})
\item[Chi-square distribution] The distribution of a sum of C squared independent standard normals; mean C, standard deviation sqrt(2C). The RX score of a Gaussian background pixel follows it. (\cref{ch:rx})
\item[Circular procedure] A procedure whose answer is needed to set up its own input; borrowing a scene's min-max frame to place a library spectrum is one. (\cref{ch:hslnorm})
\item[Cirrus] Thin, high ice cloud; EnMAP grades it in levels 0 to 3, and level 1 is treated as unknown. (\cref{ch:seghead})
\item[Class imbalance] A class covering a tiny share of pixels; per-pixel losses then let a model ignore it almost for free. (\cref{ch:seghead})
\item[Class map] Indradhanu V2's per-pixel classes: 0 land, 1 water, 2 cloud, 3 shadow, 4 snow, chosen by priority from the segmentation head's outputs. (\cref{ch:material})
\item[Class priority] The order in which classes overwrite each other in the class map: cloud over snow over shadow over water over land. (\cref{ch:shipping})
\item[Class score] The logsumexp of a class's scaled cosines; its softmax equals the sum of the per-entry softmax probabilities of the class's members. (\cref{ch:namer})
\item[Class weights] Per-class factors in the weighted mean of class losses; [1, 1, 1, 3, 3] doubles the rare classes' share of the gradient. (\cref{ch:ftruns})
\item[Class-aware residual Mahalanobis] Residual Mahalanobis run separately per class or type, each with its own whitening and a ring of same-type neighbours. (\cref{ch:scoringgains})
\item[Classical detector] An anomaly detector that computes statistics from the scene itself and has no trained weights, such as RX or MNF-RX. (\cref{ch:problem})
\item[Classifier] A model that answers with one of a fixed list of labels; it needs labelled examples and cannot return a label it was never given. (\cref{ch:problem})
\item[Cloud exclusion] Setting to 0, after scoring, the score of pixels Indradhanu V2's segmentation head calls (thick) cloud. (\cref{ch:material})
\item[Cloze test] A reading exercise in which words are blanked out and must be filled in; the analogy for masked training. (\cref{ch:masking})
\item[Codename] The Indic name (Indradhanu, Chakshu, ...) by which the backend selects a foundation model. (\cref{ch:product})
\item[Codename routing] Turning a codename into a model through a dictionary built from the catalog; a repeated codename routes to the last manifest loaded. (\cref{ch:shipping})
\item[Cofactor] The signed determinant of the block left after deleting one row and one column of a matrix; cofactors build the inverse of a 3 x 3 matrix. (\cref{ch:rx})
\item[Combined loss] Indradhanu's training objective, L1 plus lambda(t) times the mean spectral angle over graded pixels. (\cref{ch:losses})
\item[Combined score] A size-based score: 0.5 times L1 divided by its scene maximum plus 0.5 times SAM divided by its scene maximum; Indradhanu's default until 2026-10-15. (\cref{ch:inference})
\item[Common wavelength grid] The 165 wavelengths, 10 nm apart from 460 to 2450 nm with four atmospheric windows removed, onto which every hyperspectral sensor is resampled (DEFAULT\_COMMON\_WAVELENGTH\_GRID). (\cref{ch:sensors})
\item[Composite score] The weighted sum of each model's score map after min-max rescaling to [0, 1], with weights divided by their sum. (\cref{ch:material})
\item[Compressed channels (D)] The width of the spectral bottleneck; 24 in v0.1.0 and 32 from v0.2.0. (\cref{ch:compressor})
\item[computed mask] The pixels that actually received a local score, as opposed to the spatial mask of pixels eligible for one; statistics must be reduced over it. (\cref{ch:localrx})
\item[Concatenation] Joining tensors along an axis; here four 256-channel maps along the channel axis into one 1024-channel map. (\cref{ch:decoder})
\item[Condition number] The ratio of a matrix's largest to smallest eigenvalue; inverting loses roughly log10 of it in significant digits. (\cref{ch:rx})
\item[Confidence threshold] The rule that a pseudo-label counts only where the teacher's probability is above 0.9 or below 0.1. (\cref{ch:ftruns})
\item[Confounding] Two possible causes mixed into one comparison, so the result cannot be attributed to either. (\cref{ch:heavy})
\item[Confusion matrix] The two-by-two table of true positives, false positives, false negatives and true negatives at one threshold. (\cref{ch:measures})
\item[Connectivity (4 and 8)] Whether pixels that touch only at a corner count as neighbours (8) or not (4); the object table uses 8. (\cref{ch:submission})
\item[Constant extrapolation] Filling values beyond the last sample by repeating the nearest end value. (\cref{ch:bandpipeline})
\item[Context] The visible surroundings from which a hidden value is predicted. (\cref{ch:recon})
\item[Contract] What a function promises to return, provided its inputs meet stated conditions. (\cref{ch:sortbug})
\item[Contrast (hidden-pixel)] Mean error on hidden target pixels divided by mean error on hidden background pixels. (\cref{ch:decoderwork})
\item[Contrast z] Mean over bands of the absolute object-minus-ring difference, divided by the mean over bands of the ring's standard deviation. (\cref{ch:benchmarks})
\item[Control run] A run identical to the experiment except for the thing under test (FTB1: F001 with Heavy's fine-tuning recipe). (\cref{ch:heavy})
\item[Convex combination] A weighted sum whose weights are non-negative and sum to 1; it always lies between the smallest and largest of the values combined. (\cref{ch:ceiling})
\item[Convolution] A layer that slides a small kernel of shared weights over the input and takes a weighted sum at each position; in deep learning, a cross-correlation. (\cref{ch:shapes})
\item[Convolution parameter rule] A convolution has c\_in times c\_out times k squared weights plus c\_out biases. (\cref{ch:params})
\item[Coordinate reference system] The definition of what map coordinates mean, such as longitude and latitude (EPSG:4326) or metres in a UTM zone. (\cref{ch:material})
\item[Correlation] A covariance divided by both standard deviations, between -1 and 1. (\cref{ch:linalg})
\item[Cosine annealing] The schedule eta(t) = eta-min + (eta-max - eta-min)(1 + cos(pi t / T))/2: flat at the start, steepest in the middle, flat near the floor. (\cref{ch:trainloop})
\item[Cosine schedule] A learning rate that falls along half a cosine wave from its start value to a floor. (\cref{ch:heavy})
\item[Cosine similarity] The dot product of two vectors divided by the product of their lengths; for two wavelength tags it depends only on their wavelength difference. (\cref{ch:ceiling})
\item[Count channel] A channel of ones on valid pixels carried through the integral image, whose box sums give the number of valid pixels in each ring. (\cref{ch:residualmaha})
\item[Covariance] The expected product of two quantities' deviations from their means; positive when they rise and fall together. (\cref{ch:bootstrap})
\item[Covariance matrix] The average outer product of centred vectors (divided by N - 1 in the code); variances on the diagonal, pairwise covariances off it. (\cref{ch:linalg})
\item[Coverage count] The number of convolution windows that contain a given pixel; k/s on average along one axis. (\cref{ch:patchembed})
\item[Coverage pruning] Dropping bands valid in fewer than min\_valid\_pixel\_pct (default 20) percent of pixels. (\cref{ch:bandpipeline})
\item[Cross-attention] Attention in which the queries come from one set and the keys and values from another. (\cref{ch:attention})
\item[Cross-attention compressor] SpectralCrossAttentionCompressor: 32 learned queries attend over a pixel's band tokens and each answer becomes one of the 32 output channels; 5,345 parameters at d = D = 32, whatever the band count. (\cref{ch:crossstem})
\item[Cross-attention decoder] SpectralCrossAttentionDecoder: each requested wavelength becomes a query against the pixel's 32 channel tokens; 6,401 parameters at d = D = 32, for any band list. (\cref{ch:specdecoder})
\item[Cross-check (of data paths)] Scoring one model through two different reading paths on data where both exist; any gap is a bug until explained. (\cref{ch:sortbug})
\item[Cross-correlation] Sliding a kernel over an input and summing products without flipping the kernel. (\cref{ch:shapes})
\item[Cross-validation] Estimating how a setting chosen on some data will do on new data by holding data out of the choice. (\cref{ch:scoringgains})
\item[crs\_missing] The 422 error the export endpoint returns when no georeferencing can be recovered from the raw scene file. (\cref{ch:material})
\item[Cubic spline] A smooth curve of joined cubic pieces through samples; can overshoot and invent dips or bumps. (\cref{ch:bandpipeline})
\item[Curriculum] A sequence of training stages, each changing one thing and each judged before the next starts: probe, adapt, joint. (\cref{ch:finetune})
\item[Data-dependent weights] Mixing weights computed from the input itself, here a softmax of query-key scores, as opposed to the fixed weights of a convolution. (\cref{ch:crossstem})
\item[DatasetBuilder] The class that turns one sensor's raw file into a vendable. (\cref{ch:product})
\item[Datatake] One continuous EnMAP recording; its consecutive scenes border each other, so it is kept on one side of a split. (\cref{ch:pretrain})
\item[Decision table] A table of requirement, candidate, alternative and evidence, used in Part X to set HyperSL beside this product's models job by job. (\cref{ch:hslverdict})
\item[Decoder] The half of an autoencoder that rebuilds the input from the internal summary. (\cref{ch:recon})
\item[Decoder warm-up] Training only the new decoder with everything else frozen (mode decoder, FT09a) before joint training. (\cref{ch:decoderwork})
\item[Decoder width] The common channel count, 256, to which all four encoder maps are projected before they are combined. (\cref{ch:decoder})
\item[Decompressor] The bare 1 by 1 convolution from 32 compressed channels back to the band count, with no normalisation and no activation. (\cref{ch:decoder})
\item[Degrees of freedom] The number C of squared standard normals summed in a chi-square variable; for RX, the number of dimensions scored. (\cref{ch:rx})
\item[Denormalisation] The inverse map x = std times z + mean (PixelDenormalize), returning model output to reflectance. (\cref{ch:bandpipeline})
\item[Depthwise convolution] A convolution in which each channel has its own private filter and no channels are mixed (groups equal to channels); 9 weights plus a bias per channel for a 3-by-3 kernel. (\cref{ch:block})
\item[Depthwise parameter rule] A depthwise convolution over c channels has c times k squared weights plus c biases. (\cref{ch:params})
\item[Descending band list] A band list ordered from long to short wavelength, as PRISMA writes its files within each detector. (\cref{ch:sortbug})
\item[Design effect] 1 + (m - 1) rho: how much within-group correlation inflates the variance of a mean over groups of m correlated values. (\cref{ch:bootstrap})
\item[Detached tensor] A tensor cut out of the gradient computation, so a loss that uses it cannot change what produced it. (\cref{ch:uncertainty})
\item[Detector stripe] A column of pixels that all err the same way because one detector element of a push-broom spectrometer is calibrated slightly off. (\cref{ch:residualmaha})
\item[Determinant] A number computed from a square matrix (ad - bc for 2 x 2) that is zero exactly when the matrix squashes some direction flat. (\cref{ch:linalg})
\item[Diagonal loading] Adding a small multiple of the identity to a covariance before inverting (also called a ridge); every eigenvalue rises by that amount, eigenvectors unchanged. (\cref{ch:rx})
\item[Dice coefficient] Twice the overlap of two sets divided by the sum of their sizes; 1 for identical sets, 0 for disjoint ones. (\cref{ch:seghead})
\item[Difference-of-Gaussians boost] A post-filter that subtracts a wide blur from a narrow one to favour compact blobs; it caught slightly more objects but lost much AUC. (\cref{ch:residualmaha})
\item[Digital number] The raw integer a sensor file stores (0 to 65,535 for PRISMA), converted to reflectance with per-detector scale factors. (\cref{ch:pretrain})
\item[Digital number (DN)] The integer a sensor stores, proportional to detected energy; must be converted to a physical quantity before comparison. (\cref{ch:light})
\item[Dilation] Growing a binary mask by a square structuring element, optionally applied before metrics and commit. (\cref{ch:material})
\item[Directed acyclic graph (DAG)] A set of nodes joined by one-way arrows with no path that returns to its start; the shape of a Project's Actions. (\cref{ch:product})
\item[Dot product] The sum of the products of two vectors' entries, position by position; a single number, large when the vectors are long and point the same way. (\cref{ch:linalg})
\item[Drashta] Asanskrita with z-score normalisation in and out, hiding 50 to 75 percent; 267,009 parameters. (\cref{ch:convae})
\item[Drop path] Stochastic depth: dropping a whole sublayer's contribution for a whole sample; not used in this codebase. (\cref{ch:block})
\item[Dropout] During training, zeroing each value with probability p and scaling the survivors by 1/(1-p); off at evaluation. (\cref{ch:block})
\item[Dual encoder] A matching model with one network for the pixel and one for library spectra, scoring by the cosine of their embeddings, so new library entries need no retraining. (\cref{ch:material})
\item[Dual window] The square ring used by local RX: an outer window around the pixel with a smaller inner (guard) window removed; also called an annulus. (\cref{ch:localrx})
\item[Dual-window RX] The classical local detector that takes both mean and covariance from the ring; worse than the local mean on the residual. (\cref{ch:residualmaha})
\item[Early stopping] Stopping a run once its validation score has not improved by a minimum relative amount for a number of epochs called the patience. (\cref{ch:pretrain})
\item[Edge band trimming] Dropping the first and last bands (default 3) of each detector, counted on the full list of the detector's band indices. (\cref{ch:bandpipeline})
\item[Edge snapping] The patch planner's rule that moves the last patch back to end exactly at the scene edge instead of running past it. (\cref{ch:trainloop})
\item[Edge strength] In the edge diagnosis, the largest absolute difference between a pixel's band mean and that of any valid four-neighbour. (\cref{ch:ftruns})
\item[Effective mask ratio] The hidden count divided by the number of valid tokens; at most the mask ratio because of the rounding down. (\cref{ch:masking})
\item[Effective sample size] The number of independent values that would give the same variance: n m / (1 + (m - 1) rho) for n groups of m. (\cref{ch:bootstrap})
\item[Efficient self-attention (ESA)] SegFormer's attention, in which queries come from every token while keys and values come from a spatially reduced summary of the tokens. (\cref{ch:esa})
\item[Eigen-decomposition] Writing a symmetric matrix as U Lambda U transpose, with orthonormal eigenvectors in U and eigenvalues on the diagonal of Lambda. (\cref{ch:residualmaha})
\item[Eigenvalue floor] Raising any eigenvalue below 1e-12 to 1e-12 before dividing, which prevents division by zero but does not otherwise regularise. (\cref{ch:residualmaha})
\item[Eigenvector] A non-zero vector that a matrix only stretches, not turns; the stretch factor is its eigenvalue. (\cref{ch:linalg})
\item[Embedding] A vector of numbers that represents something that is not itself a vector, such as a word or a block of pixels. (\cref{ch:patchembed})
\item[Embedding dimension] The number of features per token at a stage: 32, 64, 160 and 256 for stages 1 to 4. (\cref{ch:patchembed})
\item[Embedding width] The number of features each token carries at a stage; 32, 64, 160 and 256 for stages 1 to 4. (\cref{ch:encoder})
\item[Empty token] A stage-2 token whose whole window of stage-1 tokens was hidden; about 7 of 256 at a 65 percent hiding rate. (\cref{ch:encoder})
\item[Encoder] The half of an autoencoder that turns the input into an internal summary. (\cref{ch:recon})
\item[Encoder stage] One level of the four-level SegFormer encoder: a patch embedding, two transformer blocks and a LayerNorm. (\cref{ch:patchembed})
\item[Endmember] One of the pure spectra in the linear mixing model. (\cref{ch:hslunmix})
\item[EnMAP] German hyperspectral satellite; 224 bands delivered as a folder, with ready-made cloud, cirrus, haze, shadow and snow masks. (\cref{ch:sensors})
\item[Epoch (in this project)] A fixed budget of surviving training samples (40,000 for v0.2.0) followed by validation; the interval at which the learning rate, SAM weight and checkpoints advance. Not a full pass over the data. (\cref{ch:trainloop})
\item[Epsilon (normalisation)] A small constant, 1e-5 by default in PyTorch, added to the variance before the square root so that a constant input does not cause division by zero. (\cref{ch:block})
\item[Erosion] Shrinking a validity mask so that a pixel survives only if every pixel in a square around it is valid; inference uses a 15 x 15 square. (\cref{ch:inference})
\item[Erosion radius] Half the erosion kernel, rounded down: kernel 15 removes a 7-pixel border strip, kernel 1 removes nothing. (\cref{ch:masking})
\item[Eval mode] PyTorch's model.eval(): BatchNorm uses its stored statistics and dropout is off; gradients and weight updates still work. (\cref{ch:scoringgains})
\item[Evaluation config] One JSON file naming the models, datasets, scorers, measures and groupings of an evaluation; run by scripts/evaluate.py. (\cref{ch:runeval})
\item[Evaluation mode] The mode in which BatchNorm normalises with its stored running statistics, so each output depends only on its own input. (\cref{ch:compressor})
\item[Exact ROC] The ROC built from every distinct score as a threshold, one point at the end of each run of tied scores (exact\_roc). (\cref{ch:measures})
\item[Exclude after scoring] Fitting the score's whitening on all valid pixels and then setting excluded pixels to 0, the order the benchmark measured. (\cref{ch:shipping})
\item[Exclude scorer] A wrapper that runs a base scorer and sets chosen surface classes to 0 using a segmentation head's class map. (\cref{ch:runeval})
\item[excluded\_percent] The stats.json key giving 100 times the excluded (cloud) pixels divided by the spatially valid pixels. (\cref{ch:material})
\item[Expansion ratio] The factor by which the feed-forward network widens the token in its middle; 4 here. (\cref{ch:block})
\item[Explained variance ratio] A component's eigenvalue divided by the sum of all eigenvalues: its share of the total variance. (\cref{ch:linalg})
\item[Exponential moving average] A running average in which each older value's weight shrinks by a factor beta per step. (\cref{ch:trainloop})
\item[Exponential moving average (EMA)] A running average updated as new = 0.999 old + 0.001 current; its gap to a fixed value halves in about 693 steps. (\cref{ch:ftruns})
\item[Export] The deliverable zip (GeoTIFF, Shapefile, CSV) of candidate anomalies with map coordinates. (\cref{ch:product})
\item[Export bundle] The zip produced by POST /actions/\{id\}/export: a georeferenced GeoTIFF, a Shapefile with one point per anomaly pixel, legends, a CSV and a manifest. (\cref{ch:material})
\item[Factorial design] An experiment that varies two or more factors in every combination, so each factor's main effect can be read with the others fixed. (\cref{ch:ftruns})
\item[Fallback path (ESA)] The branch taken when the token count is not H times W (tokens were removed at stage 1): no reduction, plain full attention over the surviving tokens. (\cref{ch:esa})
\item[False-alarm rate] The fraction of normal pixels a threshold flags by mistake. (\cref{ch:rx})
\item[False-positive rate (FPR)] The share of background pixels flagged at a threshold; also called the false-alarm rate. (\cref{ch:measures})
\item[Family matcher] The matcher that subtracts a ring background, scores whole families of spectra, calibrates against the scene's ordinary pixels and names a man-made family only above p*. (\cref{ch:material})
\item[FCLS] Fully constrained least squares: per pixel, the abundances that best fit the mixing model with both constraints. (\cref{ch:hslunmix})
\item[Feature map] The output of one encoder stage reshaped to (B, C, H, W): one C-number vector per grid position. (\cref{ch:encoder})
\item[Feature pyramid] A set of feature maps of the same ground at falling resolution and rising channel width; here F1 to F4 at 1/4 to 1/32 of the patch side. (\cref{ch:encoder})
\item[Feed-forward network] The per-token two-layer network of a transformer block, Linear, nonlinearity, Linear, with a wider middle. (\cref{ch:block})
\item[Fencepost formula] The number of evenly spaced points from first to last inclusive is (last - first)/step + 1. (\cref{ch:sensors})
\item[FileHelper] A class wrapping one physical file format (HE5Helper, TIFHelper, EnmapHelper, ENVIHelper, HotSatHelper) that returns raw arrays and metadata. (\cref{ch:sensors})
\item[Fill] The fraction of a pixel's area covered by the target material, between 0 and 1. (\cref{ch:namer})
\item[Fill value] A substitute number written into a dead or missing band so that a fixed-size layer can accept the pixel; it is not a measurement. (\cref{ch:ceiling})
\item[Filtered vendable] The vendable built by band\_filter\_apply on the common grid and pickled to filtered\_vendable.pkl. (\cref{ch:product})
\item[Fine-tuning] Continuing to train a pretrained model on new data or a new task, starting from its learned weights. (\cref{ch:finetune})
\item[Finite difference] An estimate of a derivative from neighbouring grid values, such as (f(i, j+1) - f(i, j-1))/2. (\cref{ch:ftruns})
\item[Fixed recipe] This book's name for a fixed linear compressor's mixing weights: the same weighted sum of bands for every pixel. (\cref{ch:ceiling})
\item[Fixed-grid model] A model that reads only the 165-band common grid (the shipped v0.2.0), so every scene must be resampled onto that grid before it reaches the model. (\cref{ch:sortbug})
\item[Flatten] An operation that merges several neighbouring axes into one without moving any number in memory. (\cref{ch:shapes})
\item[Float16] A 16-bit floating-point number; half the storage of float32 but with a largest finite value of 65,504. (\cref{ch:runeval})
\item[Float32] The 32-bit floating-point number format; its largest value is about 3.4 times 10 to the 38, so exp of anything above 88.72 overflows. (\cref{ch:attention})
\item[Floor and ceiling] The floor of u is the largest whole number not above u; the ceiling is the smallest whole number not below u. (\cref{ch:specdecoder})
\item[Foreign fraction] The share of an object's squared difference from its ring that lies outside the background subspace. (\cref{ch:benchmarks})
\item[forward\_mixed] The model method that runs one stem pass per band list, then one BatchNorm, encoder and decoder pass over the concatenated batch, and splits the output by sensor. (\cref{ch:finetune})
\item[Foundation model] In this repository, a network trained once to reconstruct ordinary imagery, whose reconstruction error is used as the anomaly score. (\cref{ch:problem})
\item[fp32 and bf16] 32-bit floating point (4 bytes per number) and 16-bit bfloat16 (2 bytes per number). (\cref{ch:crossstem})
\item[Frozen backbone] Training only a new head while every other weight, and its BatchNorm statistics, stays fixed. (\cref{ch:uncertainty})
\item[Frozen encoder] An encoder whose parameters and BatchNorm statistics are not updated during training. (\cref{ch:ftruns})
\item[Full width at half maximum] The width of a peaked response at half its height; for a Gaussian it is about 2.3548 times its standard deviation. (\cref{ch:material})
\item[Full width at half maximum (FWHM)] The width of a band's response curve measured where the response is half its peak; for a Gaussian, about 2.3548 times sigma. (\cref{ch:light})
\item[Full-attention fallback] What efficient self-attention does when tokens have been removed: it cannot rebuild the grid, so it uses all kept tokens as keys. (\cref{ch:encoder})
\item[Full-resolution skip] A second input that bypasses the encoder and reaches the decoder at full size; in FT09 it has every hidden pixel zeroed. (\cref{ch:decoderwork})
\item[Fuse convolution] The 1 by 1 convolution from 1024 to 256 channels that learns how much of each scale to use at each position. (\cref{ch:decoder})
\item[Fused attention] One GPU routine that computes scores, softmax and the weighted sum tile by tile, never storing the full weight matrix; exact, not approximate. (\cref{ch:crossstem})
\item[Gate] A pass criterion written in a run's README before launch and never moved; the next stage starts only if it is met. (\cref{ch:finetune})
\item[Gate halo] The square dilation around each gated pixel; halo 2 turns one pixel into a 5 x 5 square. (\cref{ch:decoderwork})
\item[Gated loss] A training loss that leaves out the pixels with the highest current error (FT07: the top 1\% per source group, plus a halo). (\cref{ch:decoderwork})
\item[Gather] Reading values from named positions: output k is input at index k of the index list. (\cref{ch:masking})
\item[Gating] Leaving pixels that already look anomalous out of what a network learns from. (\cref{ch:scoringgains})
\item[Gaussian cumulative distribution function] Phi(x), the probability that a standard Gaussian random number is below x. (\cref{ch:block})
\item[Gaussian distribution] The bell-shaped distribution whose density falls off as the exponential of minus half the squared number of standard deviations from the mean. (\cref{ch:rx})
\item[Gaussian mixture] A density made of K weighted Gaussian components, each modelling one background type. (\cref{ch:scoringgains})
\item[GELU] Gaussian Error Linear Unit, x times Phi(x), a smooth relative of ReLU; GELU(-1) = -0.1587, GELU(0) = 0, GELU(1) = 0.8413. (\cref{ch:block})
\item[Generalised eigenproblem] Sigma a = lambda (noise Sigma) a: find directions where one covariance is a multiple of another; solved by whitening with the second matrix and then taking ordinary eigenvectors. (\cref{ch:localrx})
\item[Geotransform (affine)] The map from a pixel's column and row to map coordinates: x = a j + b i + c, y = d j + e i + f. (\cref{ch:product})
\item[Global RX] RX with one mean and covariance estimated from all valid pixels of the scene. (\cref{ch:rx})
\item[Graded pixel] A pixel whose loss-mask value is 1, so that it contributes to the loss. (\cref{ch:losses})
\item[Gradient] The vector of partial derivatives of the loss with respect to every parameter; it points in the direction of steepest increase of the loss. (\cref{ch:trainloop})
\item[Gradient accumulation] Summing the gradients of several small batches, each loss divided by their number, before one optimiser step, to imitate a larger batch. (\cref{ch:trainloop})
\item[Gradient clipping] Scaling the whole gradient down to a maximum length when it exceeds it, keeping its direction; not used in this repository. (\cref{ch:trainloop})
\item[Gradient descent] The update that moves every parameter a small step against its gradient: w becomes w minus the learning rate times the gradient. (\cref{ch:trainloop})
\item[Gradient magnitude] The length of the vector of an image's horizontal and vertical derivatives; large on edges. (\cref{ch:ftruns})
\item[Gradient-matching loss] The edge term gradient\_l1: the L1 difference between the neighbour differences of the reconstruction and of the input, on hidden pixels. (\cref{ch:ftruns})
\item[GT-HAD] A gated-transformer hyperspectral anomaly detector trained from scratch on the scene it scores, with gating and early stopping as defences against learning the anomalies. (\cref{ch:scoringgains})
\item[Guard window] The inner square excluded from the local background, so that a target larger than one pixel does not contaminate its own statistics. (\cref{ch:localrx})
\item[HAD100] A public hyperspectral anomaly benchmark of 100 AVIRIS-NG scenes in at-sensor radiance, published with AETNet (2023). (\cref{ch:benchmarks})
\item[Half-size] The window parameter used by LocalRXDetector: a half-size of 25 means a 51 x 51 square. The residual score and PCA-RX use full widths instead. (\cref{ch:localrx})
\item[Halo] Raised scores around an anomaly caused by the anomaly contaminating its neighbours' predictions. (\cref{ch:recon})
\item[Halo ratio] The halo's contrast z over the object's; 0 for a hard edge. (\cref{ch:benchmarks})
\item[Hash] A fixed-length code computed from text that changes completely when any character changes; used to build cache keys. (\cref{ch:runeval})
\item[Hash split] Holding out a library spectrum when the first 8 hex digits of the SHA-1 of its name, over 2 to the 32 minus 1, fall below 0.2. (\cref{ch:namer})
\item[Haze colour test] Near-infrared over blue reflectance: above 1 for red bare ground, below 1 for bluish haze. (\cref{ch:runeval})
\item[Head] One of several parallel attention computations, each on its own slice of the key, query and value vectors; the stem uses 4 heads of 8 numbers. (\cref{ch:crossstem})
\item[Head dimension] The number of features each attention head works with, d = C / h. (\cref{ch:attention})
\item[Heartbeat] A timestamp a worker writes to its running job at a fixed interval to show it is alive. (\cref{ch:product})
\item[Heatmap rule] The evaluation rule that a continuous score map, or a score attribute that reads as one, disqualifies a submission. (\cref{ch:submission})
\item[Hidden-pixel rule] Train and apply a per-pixel prediction only where that pixel was hidden from the encoder, so it is judged from context alone. (\cref{ch:uncertainty})
\item[Hierarchical encoder] An encoder made of stages at successively coarser resolutions, each returning its own feature map, as opposed to a plain Vision Transformer that keeps one token grid throughout. (\cref{ch:encoder})
\item[Host] A real held-out scene that receives implanted objects. (\cref{ch:benchmarks})
\item[Hot storage] The trainer's local-disk cache of shards, synced once and reused for every epoch. (\cref{ch:trainloop})
\item[HotSat-1] Thermal sensor delivering one band of uncalibrated digital numbers around 5000. (\cref{ch:sensors})
\item[Hypernetwork] A small network whose outputs are weights for another computation; here a two-layer network maps a wavelength encoding to k, W or U. (\cref{ch:specdecoder})
\item[HyperSL] An external spectral foundation model (Kong et al., TGRS 2025): a per-pixel masked autoencoder of 9,027,841 trainable parameters, pretrained on EnMAP and DESIS L2A spectra min-max scaled per band per scene; evaluated and not adopted (Part X). (\cref{ch:hslwhat})
\item[Hyperspectral] Recording hundreds of narrow bands per pixel; high spectral, not spatial, resolution. (\cref{ch:light})
\item[Hyperspectral cube] An array of shape (C, H, W) holding one reflectance image per band; equivalently, one spectrum per pixel. (\cref{ch:light})
\item[Identity function] The function that returns its input unchanged; the perfect but useless solution to unmasked reconstruction. (\cref{ch:recon})
\item[Identity matrix] The square matrix with ones on the diagonal and zeros elsewhere; multiplying by it changes nothing. (\cref{ch:linalg})
\item[Identity transform] The geotransform with a = e = 1 and all else 0; written on every Action's GeoTIFF as a placeholder. (\cref{ch:product})
\item[Idle parameters] The 133,824 stage-1 weights (reduction convolutions, their norms and the depthwise convolutions) that receive no gradient in a masked pass because their layers are skipped. (\cref{ch:params})
\item[Implant bench] A benchmark made by adding synthetic objects of known position and difficulty to real held-out scenes. (\cref{ch:benchmarks})
\item[Implant-denoising] Training with synthetic anomalies in the input and the clean tile as the target, so the loss rewards reconstructing the background beneath an anomaly (FT05). (\cref{ch:ftruns})
\item[In-degree] The number of arrows entering a node of a directed graph. (\cref{ch:product})
\item[Inclusion-exclusion] Counting the union of two overlapping sets as the sum of their sizes minus the size of their overlap. (\cref{ch:residualmaha})
\item[Indradhanu] The hyperspectral foundation model (``rainbow''), a masked SegFormer autoencoder. (\cref{ch:problem})
\item[Indradhanu Heavy] The widened model (stage widths 96 to 768, progressive decoder) shipped as Indradhanu V2 v2.1.0; about 35.3 million parameters. (\cref{ch:params})
\item[Inference] Running a trained model on new data to obtain its predictions, with no learning taking place. (\cref{ch:inference})
\item[Inference batch size] The number of tiles passed through the network together; 8 by default in the scoring action, the main memory control on a CPU worker. (\cref{ch:inference})
\item[Information leakage] A hidden token's pixels reaching the encoder through a visible neighbour whose window overlaps them. (\cref{ch:patchembed})
\item[Input mask] In the convolutional masked models, a 0/1 channel marking pixels that are valid and visible. (\cref{ch:recon})
\item[Integral image] A table in which each entry holds the sum of every value above and to the left of it, with a leading row and column of zeros. (\cref{ch:residualmaha})
\item[Intersection over union] The size of the overlap of predicted and target sets divided by the size of their union; related to Dice by D = 2J/(1+J). (\cref{ch:seghead})
\item[Invariant 7] The evaluation rule that excluded pixels stay in the measures as negatives tied at 0, and labelled objects inside them stay as unfindable positives. (\cref{ch:runeval})
\item[Inverse] The matrix that undoes a square matrix, so that their product is the identity; it exists only when the determinant is non-zero. (\cref{ch:linalg})
\item[Inverse map] A map that undoes another: D(N(x)) = x and N(D(z)) = z. (\cref{ch:bandpipeline})
\item[Inverse permutation] The permutation that undoes another: it gives, for each token, the place it occupies in the sorted order (its rank). Called ids\_restore in the original masked autoencoder. (\cref{ch:masking})
\item[Irradiance] Power of light arriving on a surface per square metre per nanometre. (\cref{ch:light})
\item[Jacobian] The matrix of partial derivatives of a vector function with respect to its vector input; the vector form of a derivative. (\cref{ch:block})
\item[Joint stage] Training reconstruction and segmentation together: loss = reconstruction + seg weight x segmentation. (\cref{ch:finetune})
\item[Jump] The distance in input pixels between neighbouring outputs of a layer; the product of the strides so far. (\cref{ch:shapes})
\item[k-means] Lloyd's clustering recipe: assign each point to its nearest centre, move each centre to the mean of its points, repeat; the hard-assignment limit of a mixture. (\cref{ch:scoringgains})
\item[Keep mask] The per-token mask saying which stage-1 tokens are shown to the encoder; when it removes tokens, stage 1 attention takes the fallback path. (\cref{ch:esa})
\item[Kernel] The small block of weights a convolution slides over its input; its side length is the kernel size k. (\cref{ch:shapes})
\item[Kernel span] The wavelength range covered by a convolution window over band index: the last band's wavelength minus the first's. (\cref{ch:hslnorm})
\item[Key] The learned projection of a token that expresses what it advertises about itself; compared against queries. (\cref{ch:attention})
\item[Key count] The number of keys efficient self-attention uses, N divided by R squared; 16 at every stage of a 128-pixel patch. (\cref{ch:encoder})
\item[kmat] The matrix of shape (bands, bins x 32) that places each band's k(lambda) in its bin's slot, so one matrix product forms all bin tokens without a per-band tensor. (\cref{ch:specdecoder})
\item[Known mask] A per-class, per-pixel flag saying whether a label is trusted; unknown pixels are left out of the loss and the IoU. (\cref{ch:seghead})
\item[L1 loss] The average absolute error; its best constant guess is the median, so it is robust to rare extreme values. (\cref{ch:recon})
\item[L1 norm] The sum of the absolute values of a vector's entries. (\cref{ch:losses})
\item[L2 loss] The average squared error; its best constant guess is the mean, so rare extreme values pull predictions towards them. (\cref{ch:recon})
\item[L2 regularisation] Adding half lambda times the squared length of the weights to the loss; equal to weight decay under plain gradient descent but not under Adam. (\cref{ch:trainloop})
\item[L2-SP] A regulariser that adds alpha/2 times the squared distance of the weights from their pretrained values, pulling them back during fine-tuning. (\cref{ch:ftruns})
\item[Label] The correct answer attached to a training example. (\cref{ch:problem})
\item[Label audit] Checking labels without ground truth by testing whether what is labelled looks physically like its class, and finding what predicts failures. (\cref{ch:seghead})
\item[Label ceiling] The limit on a class's accuracy set by wrong labels: a class that stays flat while every other class improves. (\cref{ch:ftruns})
\item[Label policy] The fixed rules (no-data, cloud class, cirrus level 1, overlaps, dry scenes) that decide which EnMAP labels are known; set from the audit before training. (\cref{ch:seghead})
\item[Lagrange multiplier] A number introduced to maximise a function under a constraint; at the optimum the function's gradient is that multiple of the constraint's gradient. (\cref{ch:linalg})
\item[Lambertian surface] An idealised matt surface that looks equally bright from every direction, for which total exitance equals pi times radiance. (\cref{ch:light})
\item[Land surface temperature] The temperature of the ground's skin inferred from the heat it radiates; Chakshu's input, in degrees Celsius for Landsat 9. (\cref{ch:chakshu})
\item[Landsat 9] Thermal source in this product: one band (B10) delivered as a GeoTIFF, converted to degrees Celsius. (\cref{ch:sensors})
\item[Laplace distribution] The probability distribution with density (1/2b) exp(-|r|/b); sharp-peaked at zero with exponential tails, and with mean absolute value equal to its scale b. (\cref{ch:uncertainty})
\item[Last-wins] The rule that assigning to an existing dictionary key, or to an array element, overwrites the earlier value. (\cref{ch:shipping})
\item[Law of total variance] Total variance equals the mean of the within-class variances plus the variance of the class means. (\cref{ch:finetune})
\item[Layer normalisation (LayerNorm)] Standardising each token over its own features, then applying a learned scale and shift per feature; independent of the batch. (\cref{ch:patchembed})
\item[LayerNorm] Normalisation of one token across its own C features: subtract their mean, divide by their standard deviation, then apply a learned per-feature scale and shift; 2C parameters. (\cref{ch:block})
\item[LayerNorm parameters] One learned scale and one learned shift per feature: 2c numbers. (\cref{ch:params})
\item[Layout] The order in which a multi-index array is flattened into memory. (\cref{ch:sensors})
\item[Lazy-import rule] Action modules may import app, torch or rasterio only inside functions, keeping the api's start-up light. (\cref{ch:product})
\item[Leakage] A benchmark scene overlapping or neighbouring training data, so a model can score by memory. (\cref{ch:benchmarks})
\item[Learned query] A query vector stored as a trained parameter rather than computed from the input; the spectral compressor has 32 of them. (\cref{ch:ceiling})
\item[Learning rate] The step-size multiplier eta in a gradient update; too small crawls, too large diverges. (\cref{ch:trainloop})
\item[Learning-rate schedule] A rule that changes the learning rate over the course of training. (\cref{ch:trainloop})
\item[Leave-one-scene-out] Cross-validation with scenes as units: each scene is reported at the setting chosen on the other scenes only. (\cref{ch:scoringgains})
\item[Leave-one-spectrum-out] Matching a planted library spectrum against every other library spectrum, never itself, so the test asks whether the nearest other spectrum shares its family. (\cref{ch:material})
\item[Level 1 by matching] The namer's natural-or-man-made answer since run N001b: the natural library entries pooled as an eighth class, and P(man-made) = sigmoid(LSE of the seven man-made class scores - the natural score + b), the share of softmax mass on man-made entries. (\cref{ch:namer})
\item[Lexicographic order] Sorting strings by comparing character codes from the left; a (97) before h (104) before s (115). (\cref{ch:shipping})
\item[Likelihood] The probability a model gives to the observed labels, viewed as a function of the model's predictions. (\cref{ch:seghead})
\item[Likelihood ratio] The probability density of the data under one hypothesis divided by that under another; RX is the generalised likelihood ratio for an unknown shift from a Gaussian background. (\cref{ch:rx})
\item[Linear bottleneck] A pair of linear maps through a narrow width; every output lies in a subspace whose dimension is at most that width. (\cref{ch:compressor})
\item[Linear interpolation] Estimating a value between two samples on the straight line through them: y = (1 - t) y0 + t y1 with t the fraction of the way. (\cref{ch:bandpipeline})
\item[Linear interpolation (percentiles)] np.percentile's rule: the q-th percentile of n sorted values sits at position (n - 1) q / 100, interpolated between neighbours. (\cref{ch:bootstrap})
\item[Linear layer] A layer computing y = W x + b; it has in times out plus out parameters. (\cref{ch:shapes})
\item[Linear map] A map f with f(ax + cy) = a f(x) + c f(y); its sensitivity to each input is a constant, the same for every input vector. (\cref{ch:ceiling})
\item[Linear mixing model] A pixel covering fraction a of material t and 1 - a of background b reads (1 - a) b + a t. (\cref{ch:benchmarks})
\item[Linear mixture] A synthetic pixel made as a fraction of a target spectrum plus the rest of the background spectrum. (\cref{ch:ftruns})
\item[Linear-interpolation percentile] numpy's default percentile: the value at position q/100 times (N - 1) in the sorted list, interpolated between neighbours. (\cref{ch:material})
\item[Linearity of expectation] The expected value of a sum equals the sum of the expected values, whether or not the terms are independent. (\cref{ch:encoder})
\item[Loader worker] A separate process that reads and decodes data while the GPU trains; each needs at least one shard. (\cref{ch:heavy})
\item[Local mean] The average of a quantity over a ring of neighbouring pixels, used in place of the scene-wide mean. (\cref{ch:residualmaha})
\item[Local RX] RX with the mean and covariance estimated separately for each pixel from a ring of its neighbours, so a pixel is judged against its surroundings rather than the whole scene. (\cref{ch:localrx})
\item[Local z-score] A spatial post-filter that re-expresses each score against its neighbourhood's mean and spread; it lowered AUC in the comparison. (\cref{ch:residualmaha})
\item[Local-mean Mahalanobis distance] The Mahalanobis distance measured from the ring mean but with the scene-wide covariance; the default with a 21/61 window. (\cref{ch:residualmaha})
\item[Log parameterisation] Letting a network output s = log b and using b = exp(s), so b is always positive and gradients near zero stay moderate. (\cref{ch:uncertainty})
\item[Log-odds] The logarithm of p/(1-p); it runs over all real numbers as p runs from 0 to 1. (\cref{ch:seghead})
\item[Log-uniform distribution] A distribution whose logarithm is uniform; the namer's fill on [0.05, 1], with P(a <= q) = ln(q / 0.05) / ln 20 and mean 0.95 / ln 20 = 0.3171. (\cref{ch:namer})
\item[Logit] A raw real-valued score, read as the log-odds of yes; the sigmoid turns it into a probability. (\cref{ch:seghead})
\item[Logsumexp] ln of the sum of exponentials, a smooth maximum; the namer's class score over a class's library cosines divided by tau. (\cref{ch:namer})
\item[Loss] A single number measuring how bad a model's output is; training lowers it. (\cref{ch:recon})
\item[Loss mask] The pixel prediction mask times the eroded validity mask: the pixels the loss is averaged over. (\cref{ch:masking})
\item[Machine epsilon] The relative rounding step of a floating-point format: about 1.2e-7 for float32 and 2.2e-16 for float64. (\cref{ch:rx})
\item[Machine epsilon (float32)] The relative spacing of float32 numbers near 1, 2 to the power -24 below 1; it makes cosines of angles under about 0.00024 rad round to exactly 1. (\cref{ch:losses})
\item[Mahalanobis distance] (x - mean)T inverse-covariance (x - mean): the squared distance from the mean after whitening, so each direction is measured in its own standard deviations. (\cref{ch:rx})
\item[Manhattan distance] Rows apart plus columns apart; the number of 4-neighbour steps between two pixels. (\cref{ch:benchmarks})
\item[Manifest (current.json)] The file beside a model's weights naming which checkpoint is current and recording facts about it. (\cref{ch:product})
\item[Manifest parameter figure] The count recorded with a checkpoint; it includes buffers, so it exceeds the trainable count (5,507,354 against 5,506,629). (\cref{ch:params})
\item[Mann-Whitney count] The number of (target, background) pairs in which the target scores higher, ties counting half; divided by the number of pairs it equals the AUC. (\cref{ch:measures})
\item[Mask] An array of zeros and ones selecting positions; in masked reconstruction 1 means visible and 0 means hidden. (\cref{ch:recon})
\item[Mask ratio] The fraction of tokens hidden in a training step; 0.65 in Indradhanu's v0.2.0 training configuration. (\cref{ch:recon})
\item[Mask token] In the original masked autoencoder, a learned vector placed at hidden positions; this encoder uses zeros instead. (\cref{ch:encoder})
\item[Masked autoencoder (MAE)] An autoencoder that hides part of its input and is graded only on the hidden part, so it cannot copy. (\cref{ch:recon})
\item[Masked mean] A mean over padded rows that sums only the real entries and divides by their count, not by the padded length. (\cref{ch:localrx})
\item[Masked mean absolute error (L1 loss)] The mean of the absolute errors over all bands of all graded pixels; Indradhanu's brightness term, in reflectance units. (\cref{ch:losses})
\item[Masked softmax] A softmax in which invalid positions have their scores set to minus infinity, giving them weight exactly 0 and renormalising the rest. (\cref{ch:crossstem})
\item[Matched filter] The linear filter w = Sigma-inverse t / (t-transpose Sigma-inverse t) that maximises signal-to-noise for a known signal shape t in noise of covariance Sigma. (\cref{ch:scoringgains})
\item[Matrix] A rectangular table of numbers; multiplying a vector by it forms one weighted mixture of the inputs per row. (\cref{ch:linalg})
\item[Max invalid voxel fraction] The stage-6 threshold (default 0.4): a pixel invalid in more than this fraction of the surviving bands is invalidated in every band. (\cref{ch:bandpipeline})
\item[Max pooling] Sliding a window over a grid and outputting the maximum of each window; on a 0/1 mask it asks whether any 1 is in reach. (\cref{ch:masking})
\item[Max pooling as dilation] Taking the maximum of a binary mask over a sliding window, which grows every marked pixel into a square. (\cref{ch:decoderwork})
\item[Max-normalisation] Dividing a map by its largest valid value; fragile, because one extreme pixel compresses all the others towards zero. (\cref{ch:inference})
\item[Maximum likelihood] Choosing the parameter value that makes the observed data most probable. (\cref{ch:uncertainty})
\item[Maximum noise fraction] The MNF transform (Green et al. 1988): the directions ordered by noise fraction, equivalently by signal-to-noise ratio; often called minimum noise fraction. (\cref{ch:localrx})
\item[max\_kept leakage] When patches in a batch keep different numbers of tokens, every patch gathers the batch maximum, so patches that keep fewer let a few hidden tokens into the encoder. (\cref{ch:masking})
\item[McClish standardisation] A rescaling of the raw partial area that maps random to 0.5 and perfect to 1; not used by this repository. (\cref{ch:measures})
\item[Mean] The sum of a set of numbers divided by how many there are. (\cref{ch:problem})
\item[Mean squared error (MSE)] The mean of the squared errors; its gradient is proportional to the error, so large errors dominate it. (\cref{ch:losses})
\item[Mean vector] The band-by-band average of a set of spectra. (\cref{ch:linalg})
\item[Mebibyte (MiB)] 2 to the power 20 = 1,048,576 bytes; a gibibyte (GiB) is 1,073,741,824 bytes. (\cref{ch:crossstem})
\item[Median] The middle value of sorted observations; the constant prediction that minimises the sum of absolute errors. (\cref{ch:losses})
\item[Memory per extra patch] The slope of peak memory against batch size; it removes fixed costs such as weights and optimiser state from a per-patch comparison. (\cref{ch:specdecoder})
\item[Memory wall] The point at which memory, not arithmetic or parameters, limits the batch size and speed of training. (\cref{ch:crossstem})
\item[MiB] A mebibyte, 2 to the 20 = 1,048,576 bytes; a float32 number takes 4 bytes. (\cref{ch:esa})
\item[Micro-batch] The batch of one forward pass (8 tiles here); BatchNorm statistics are computed per micro-batch. (\cref{ch:finetune})
\item[min bg pixels] The least number of valid ring pixels needed to score a pixel locally, by default the number of bands plus one, because a covariance from n points has rank at most n - 1. (\cref{ch:localrx})
\item[Min-max rescaling] Mapping a map's values to [0, 1] by subtracting its minimum and dividing by its range. (\cref{ch:material})
\item[Minimum Noise Fraction (MNF)] A classical linear projection to few bands chosen by eigen-decomposition to maximise signal-to-noise ratio. (\cref{ch:compressor})
\item[Mix-FFN] SegFormer's feed-forward network with a 3-by-3 depthwise convolution on the token grid between the expansion and the GELU. (\cref{ch:block})
\item[Mixdown] The analogy for compression: each output channel is a learned blend of every input band. (\cref{ch:compressor})
\item[Mixed background] A scene made of several distinct surfaces, whose global mean and covariance describe none of them well. (\cref{ch:rx})
\item[Mixed batch] A micro-batch holding tiles from every sensor, each sensor's sub-batch passed through the stem alone and then concatenated. (\cref{ch:finetune})
\item[Mixed pixel] A pixel whose spectrum is a combination of several materials because each covers part of the ground it sees. (\cref{ch:material})
\item[Mixed precision] Running the expensive parts of a network in a 16-bit floating-point format and the precision-critical parts, such as the loss, in float32. (\cref{ch:trainloop})
\item[Mixing line] The straight segment in spectral space from the background spectrum to the material spectrum, on which every two-part mixture lies. (\cref{ch:benchmarks})
\item[MNF component] A pixel's coordinate along an MNF direction; MNF components are uncorrelated, with variance 1 + SNR. (\cref{ch:localrx})
\item[MNF eigenvalue diagnostic] Reading the MNF eigenvalues: values well above 1 are signal components, values near 1 are noise; a flat list means no dominant signal. (\cref{ch:localrx})
\item[MNF-RX] RX computed after a Minimum Noise Fraction transform into a smaller, better-conditioned space; the shipped classical hyperspectral detector. (\cref{ch:problem})
\item[Mode scale] The fine-tuning trainer mode that trains only the uncertainty head with the Laplace NLL and ranks epochs by that NLL. (\cref{ch:uncertainty})
\item[Model-free baseline] A detector with no learned model, run beside the models to check that a benchmark rewards what a model adds. (\cref{ch:runeval})
\item[Momentum] An optimiser that steps along a running sum of past gradients, so consistent directions accumulate and alternating ones cancel. (\cref{ch:trainloop})
\item[Momentum (BatchNorm)] The weight m given to the newest batch statistic when updating the running statistics; 0.1 by default. (\cref{ch:compressor})
\item[mRMSE] Mean root-mean-square error of the abundance maps after matching. (\cref{ch:hslunmix})
\item[mSAD] Mean spectral angle, in radians, between matched found and true endmembers. (\cref{ch:hslunmix})
\item[Multi-head attention] Several attentions run side by side on separate slices of the features, concatenated and mixed by an output projection. (\cref{ch:attention})
\item[Multi-label] A classification problem in which any number of classes can be true at once, such as thin cloud over water. (\cref{ch:seghead})
\item[Multiply-add] One multiplication followed by one addition; the unit in which convolution work is counted. (\cref{ch:decoder})
\item[Multiply-add (MAC)] One multiplication followed by one addition, the unit used here to count the work of matrix products and convolutions. (\cref{ch:esa})
\item[Multiply-adds] The work of a layer: its weights per output position times the number of positions. Not the same as its parameter count. (\cref{ch:params})
\item[Multivariate Gaussian] The Gaussian over vectors, set by a mean vector and covariance matrix; its exponent is minus half the Mahalanobis distance. (\cref{ch:rx})
\item[NaN] Not a number: the floating-point result of 0/0, produced by a softmax over all-masked scores unless guarded against. (\cref{ch:crossstem})
\item[Nanometre] One billionth of a metre (10 to the power -9 m); the unit of wavelength used throughout the codebase. (\cref{ch:light})
\item[Native band list] A sensor's own band centres and pixel grid, untouched by interpolation. (\cref{ch:pretrain})
\item[Native bands] The wavelengths a sensor actually measured, as opposed to the 165-band common grid; the band-filter default since 2026-10-31. (\cref{ch:material})
\item[NDSI] Normalised difference snow index, (green - SWIR)/(green + SWIR); high for snow, which is dark near 1600 nm where cloud stays bright. (\cref{ch:seghead})
\item[NDVI] Normalised difference vegetation index, (NIR - red) / (NIR + red) or with a visible mean; positive for vegetation, negative for water. (\cref{ch:runeval})
\item[Near infrared (NIR)] Wavelengths from about 700 to 1000 nm, immediately beyond red; vegetation reflects strongly here. (\cref{ch:light})
\item[Near weight] The factor (5 in FT05) applied to hidden pixels within 3 px of an implant in the loss mask; it changes shares of a weighted mean, not its size. (\cref{ch:ftruns})
\item[Nearest-valid fill] Replacing each invalid pixel's values with those of the nearest valid pixel by Euclidean distance, leaving the validity masks unchanged. (\cref{ch:bandpipeline})
\item[needs\_threshold] The status in which anomaly\_detection\_prep waits for a person to choose and commit a threshold. (\cref{ch:product})
\item[Negative (background)] A valid pixel the ground truth does not mark as anomalous. (\cref{ch:measures})
\item[Negative log-likelihood] Minus the logarithm of the likelihood; minimising it is the same as maximising the likelihood, and it turns products over pixels into sums. (\cref{ch:uncertainty})
\item[Neighbouring checkpoints] Consecutive epochs' weights; a gain counts only if it holds at all three compared. (\cref{ch:heavy})
\item[Neural network] A function built from many simple layers whose behaviour is set by a large list of learned numbers called parameters. (\cref{ch:problem})
\item[Neuron] A weighted sum of inputs plus a bias, usually followed by an activation function. (\cref{ch:shapes})
\item[Nodata] The value -32768 in a shard, marking a voxel with no valid reading; the stem skips it rather than reading it as zero. (\cref{ch:pretrain})
\item[Noise fraction] Along a direction, the share of the total variance that is noise. (\cref{ch:localrx})
\item[Noise whitening] Transforming data so the noise covariance becomes the identity, the first step of MNF. (\cref{ch:localrx})
\item[Non-overlapping windows] Convolution windows with kernel equal to stride and no padding, so each pixel feeds exactly one token; Indradhanu's stage 1. (\cref{ch:patchembed})
\item[Norm] The length of a vector, the square root of the sum of its squared entries (Pythagoras in many dimensions). For a spectrum, roughly its brightness. (\cref{ch:linalg})
\item[Normalised convolution] Blurring data with nodata set to 0 and dividing by the blurred validity mask, so missing pixels do not drag values down. (\cref{ch:scoringgains})
\item[np.allclose] NumPy's test that two arrays agree within atol + rtol times the size of each value (defaults 1e-8 and 1e-5). (\cref{ch:shipping})
\item[np.interp] NumPy's one-dimensional linear interpolator, np.interp(x, xp, fp); it requires the sample positions xp to be increasing and does not check. (\cref{ch:sortbug})
\item[Null space] The set of input directions a matrix sends to zero; changes along them are invisible after the matrix. (\cref{ch:compressor})
\item[Numerical guard] A clamp or small added constant (epsilon) that keeps a computation away from a division by zero, a negative square root or an infinite derivative. (\cref{ch:losses})
\item[Object] A 4-connected component of the ground-truth mask. (\cref{ch:measures})
\item[Object AUC] The share of background pixels an object's best-scoring pixel outranks. (\cref{ch:hslunmix})
\item[Object bootstrap] Resampling labelled objects rather than pixels, because an object's pixels are not independent evidence. (\cref{ch:bootstrap})
\item[Object measure] A detection measure that counts each object once, such as objects caught at the top 1 percent. (\cref{ch:scoringgains})
\item[Object pAUC parts] The pair (sum of c over an object's pixels, its pixel count) from which a bootstrap rebuilds pAUC. (\cref{ch:bootstrap})
\item[Objects caught] The share of objects with a pixel within one pixel of the top q percent of valid scores. (\cref{ch:measures})
\item[Onboarding vendable] The vendable built at upload, with the sensor's native bands. (\cref{ch:product})
\item[Open-set recognition] Recognition that may meet classes it never saw in training and say so; here, level 1 calls man-made whatever its model of nature cannot explain. (\cref{ch:material})
\item[Opt-in model] A model available to choose but not used by any default. (\cref{ch:shipping})
\item[Ordering test] Acceptance of a benchmark only if it reproduces a ranking the real labels already showed: fine-tuned models above a model-free PCA-RX. (\cref{ch:benchmarks})
\item[Orthogonal] At right angles; two vectors are orthogonal when their dot product is zero. (\cref{ch:linalg})
\item[Orthogonal matrix] A square matrix whose columns are orthogonal unit vectors; it rotates or reflects without stretching, and its inverse is its transpose. (\cref{ch:linalg})
\item[Orthogonal projection] The closest point of a span to a vector; the leftover is at right angles to the span. (\cref{ch:decoderwork})
\item[Outer product] The matrix d dT whose entry (i, j) is the product of entries i and j; averaging outer products of deviations gives the covariance. (\cref{ch:linalg})
\item[Output projection] The final C-by-C linear layer of an attention module, which lets the separate heads' outputs interact. (\cref{ch:attention})
\item[Output-size formula] For kernel k, stride s and padding p, an axis of length H becomes floor((H + 2p - k)/s) + 1. (\cref{ch:shapes})
\item[Overall accuracy] The share of pixels whose predicted (or cluster-mapped) class equals their true class. (\cref{ch:hslwhat})
\item[Overfitting] Learning the training examples rather than the pattern: training loss falls while validation loss rises. (\cref{ch:heavy})
\item[Overlap] The number of pixels two neighbouring tiles share along an axis, patch size minus stride: 96 for a 128 tile at stride 32. (\cref{ch:inference})
\item[Overlap averaging] Combining the predictions of all tiles that cover a pixel by summing them, weighted by the eroded validity mask, and dividing by the sum of weights. (\cref{ch:inference})
\item[OverlapPatchEmbedding] The repository's patch-embedding module: Conv2d, flatten, transpose, LayerNorm; returns the tokens and the grid height and width. (\cref{ch:patchembed})
\item[Overlapping windows] Convolution windows with kernel larger than stride, so neighbouring tokens share pixels; Indradhanu's stages 2 to 4. (\cref{ch:patchembed})
\item[Overshoot] An interpolated curve passing below the lowest or above the highest nearby sample. (\cref{ch:bandpipeline})
\item[P(B>A)] The share of bootstrap draws in which scorer B's measure beat scorer A's. (\cref{ch:bootstrap})
\item[p*] The percentile bar 1 - alpha / F for giving a man-made name, chosen so that an ordinary pixel gets any man-made name with probability at most alpha. (\cref{ch:material})
\item[Padding] A border of zeros added around the input before a convolution. (\cref{ch:shapes})
\item[Paired bootstrap] A bootstrap that applies the same draw to both scorers being compared, so that which objects were drawn cancels in the difference. (\cref{ch:bootstrap})
\item[Parameter] One adjustable number inside a model, set by training. (\cref{ch:problem})
\item[Parameter count] The total number of parameters; 5,506,629 for the 165-band Indradhanu v0.2.0. (\cref{ch:params})
\item[Parity run] Running the shipped checkpoint through the product's own code path on benchmark scenes and comparing with the evaluator's numbers. (\cref{ch:shipping})
\item[Partial AUC (pAUC)] The area under the ROC from FPR 0 to a limit, here divided by the limit so that perfect is 1 and random is half the limit (0.005 at 1\%). (\cref{ch:measures})
\item[Partition] A split of a set into pieces that do not overlap and together cover the whole set; at inference, the hidden token sets of the passes form a partition. (\cref{ch:inference})
\item[Pass criterion (gate)] A condition a run must meet, written in its README before launch and not moved afterwards. (\cref{ch:ftruns})
\item[Patch] A small square cut from a scene, here 128 by 128 pixels with every band, the unit a training batch is built from. (\cref{ch:pretrain})
\item[Patch embedding] The operation that cuts a feature map into blocks and computes one vector, a token, per block; in this model a strided convolution followed by LayerNorm. (\cref{ch:patchembed})
\item[Patience] The number of epochs without sufficient improvement after which early stopping ends a run. (\cref{ch:pretrain})
\item[PCA-RX] The evaluation's classical baseline: whiten the top 30 principal components of the cube with a global covariance and score the squared distance from a ring mean (full widths 21 and 61). (\cref{ch:localrx})
\item[PCHIP] Piecewise cubic Hermite interpolating polynomial: cubic pieces with shape-preserving slopes, never creating an extremum absent from the data; used for stage-7 gap filling. (\cref{ch:bandpipeline})
\item[Per-band affine map] A map that multiplies each band by its own number and adds its own number; min-max and the fixed z-score are both of this form. (\cref{ch:hslnorm})
\item[Per-band normalisation] Rescaling each band by its own corpus mean and standard deviation (PixelNormalize), done inside the model. (\cref{ch:bandpipeline})
\item[Per-curve scaling] Rescaling one spectrum by its own minimum and maximum over its bands; used in the HyperSL matching run so that pixel and library were treated alike. (\cref{ch:hslnorm})
\item[Per-pixel contribution c(p)] max(0, m - FPR(p)) / m, with FPR(p) the share of background scoring at or above target pixel p; its mean over target pixels is the pAUC. (\cref{ch:bootstrap})
\item[Per-pixel gate] SpectralGate: a small network computes one sigmoid volume per band from the pixel's own spectrum before the fixed 1x1 compressor; 13,162 parameters at 202 bands. (\cref{ch:specdecoder})
\item[Per-scene min-max normalisation] Rescaling each scene by its own minimum and maximum; avoided because it destroys absolute reflectance. (\cref{ch:specdecoder})
\item[Per-scene min-max scaling] Rescaling each band so that the scene's darkest pixel in that band is 0 and its brightest 1: x' = (x - min) / (max - min), with min and max taken over the scene. HyperSL's input scaling. (\cref{ch:hslnorm})
\item[Per-scene override] Replacing a model's stored normalisation mean and spread with the scene's own (PixelStatsOverride); meant for HotSat-1, but never built, and overwritten by the checkpoint if it were. (\cref{ch:chakshu})
\item[Percentile interval] The 2.5th to 97.5th percentiles of the bootstrap values: a 95 percent interval. (\cref{ch:bootstrap})
\item[Percentile threshold] A cut that keeps a fixed fraction of a scene's highest scores, so it adapts to the scene's own score scale. (\cref{ch:problem})
\item[Period] The distance, in nm, after which a sine or cosine repeats; the encoding uses five periods per decade. (\cref{ch:ceiling})
\item[Permutation] A reordering of the indices 0 to N-1 in which each index appears exactly once. (\cref{ch:masking})
\item[Permutation equivariance] The property that reordering the input tokens reorders the outputs the same way and changes nothing else; plain attention has it, so it cannot tell where a token is. (\cref{ch:attention})
\item[Permutation invariance] The property that reordering the keys and values (together) does not change an attention output. (\cref{ch:ceiling})
\item[Pickle] Python's object serialisation; stores the class's import path, which makes class names a wire format. (\cref{ch:product})
\item[Pixel] One cell of an image grid, standing for a patch of ground; named by its row and column, counted from 0. (\cref{ch:problem})
\item[Pixel centre] The middle of a pixel, half a pixel in from its edges; PRISMA's corner attributes give the corner pixels' centres. (\cref{ch:submission})
\item[Pixel checkerboard] The inference mask of the convolutional models: one pass hides pixels whose row plus column is odd, the other those where it is even. (\cref{ch:convae})
\item[Pixel measure] A detection measure that counts target pixels, such as pAUC or TPR; a large object weighs more than a one-pixel target. (\cref{ch:scoringgains})
\item[Pixel prediction mask] The token prediction mask copied to each token's 4 by 4 pixels by nearest-neighbour upsampling. (\cref{ch:masking})
\item[Pixel shuffle] A parameter-free rearrangement from (B, C r squared, H, W) to (B, C, rH, rW): out[c, rh + i, rw + j] = in[c r squared + r i + j, h, w]. (\cref{ch:decoder})
\item[Pixel validity mask] A (B, 1, H, W) array of ones and zeros: 1 where the pixel is a real measurement, 0 for cloud, no-data or swath edge. (\cref{ch:masking})
\item[Pixel with no valid band] A pixel whose every band is masked; the code lets it attend to everything and then sets its output to exactly 0. (\cref{ch:crossstem})
\item[Pixel-shuffle bottleneck] The decoder's step that turns one quarter-resolution feature vector into a 4 x 4 block of pixels, so sub-block detail must be carried in that one vector. (\cref{ch:ftruns})
\item[Pixels as a batch] Treating each pixel of a patch as an independent row (16,384 rows per 128 x 128 patch) because nothing in the spectral stem moves information between pixels. (\cref{ch:crossstem})
\item[PixelStatsOverride] Per-scene replacement normalisation statistics for sensors whose units differ from the model's training data; makes scores scene-relative. (\cref{ch:bandpipeline})
\item[Placement bias] A benchmark bias in which synthetic targets are placed on surfaces where real targets do not occur, so a scorer can win on the benchmark and lose on real scenes. (\cref{ch:scoringgains})
\item[Plain RX] The Mahalanobis distance of each pixel of the raw cube over all bands; dropped from the product in 2026-05 because its near-singular covariance gave distances near 1e11. (\cref{ch:residualmaha})
\item[Planted-sample generator] The namer's source of training data: x = g (a t + (1 - a) b) + n, a library spectrum mixed into a real pixel with a gain curve and noise. (\cref{ch:namer})
\item[Plateau] A stretch of training in which further epochs no longer change the measured result beyond noise. (\cref{ch:heavy})
\item[Plausible cloud] A labelled cloud pixel that is bright (visible mean above 0.2) and flat (min/max of four window means at least 0.7). (\cref{ch:seghead})
\item[Point-spread function (PSF)] The image of a single bright point; here a Gaussian of standard deviation 1.4 pixels. (\cref{ch:benchmarks})
\item[Position code] A fixed vector added to each token to say where it sits; HyperSL's is sin and cos of (wavelength - 400) / 10 at 128 frequencies 100 to the power -2j/256. (\cref{ch:hslwhat})
\item[Positional encoding] A vector added to each token to tell a transformer where the token sits; Indradhanu has none and gets position from a depthwise convolution instead. (\cref{ch:attention})
\item[Positive (target)] A valid pixel the ground truth marks as anomalous. (\cref{ch:measures})
\item[Positive semi-definite] A symmetric matrix with no negative eigenvalues; every covariance matrix is one, because its eigenvalues are variances. (\cref{ch:linalg})
\item[Post-norm] Placing LayerNorm after the residual addition, LN(x + f(x)), as in the original 2017 transformer. (\cref{ch:block})
\item[Pratibimba] The plain convolutional autoencoder (reflection): four stages, base 32, trained with nothing hidden; 1,379,393 parameters. (\cref{ch:convae})
\item[Pre-norm] Placing LayerNorm inside the residual branch, x + f(LN(x)), so that the residual path is never normalised. (\cref{ch:block})
\item[Precision] For a class, the share of the pixels the classifier calls it that are labelled with it: diagonal entry over column total. (\cref{ch:hslcloud})
\item[Precondition] A condition the inputs of a function must satisfy for its contract to hold; when it is violated, the function promises nothing. (\cref{ch:sortbug})
\item[Prediction mask] The (B, N) array of hidden tokens, the ones the model is graded on. Only valid tokens can be in it. (\cref{ch:masking})
\item[Prep bundle] The export of a committed anomaly\_detection\_prep: a 0/1 mask GeoTIFF, a point shapefile without a score attribute, a CSV and a manifest, named hyper\_ or thermal\_ by sensor; the only bundle fit for submission. (\cref{ch:submission})
\item[Pretraining] Training a model on a large general data set before adapting it (fine-tuning) to the data and tasks of the product. (\cref{ch:pretrain})
\item[Principal component analysis] Finding the orthogonal directions of greatest variance of a cloud of vectors: the eigenvectors of its covariance, ordered by eigenvalue. (\cref{ch:linalg})
\item[PRISMA] Italian Space Agency hyperspectral satellite; 239 bands (66 VNIR, 173 SWIR) delivered as one HDF5 (.he5) file in BIL layout. (\cref{ch:sensors})
\item[Probability density] A curve (or surface) whose area (or volume) over a region is the probability of landing in that region. (\cref{ch:rx})
\item[Probe] Training only a small head on frozen features, to measure what the features already know. (\cref{ch:finetune})
\item[Product grid (PRISMA)] The UTM grid of a PRISMA L2D product, read from its EPSG code and the corner pixels' centres: pitch = span / (n - 1), origin = upper-left centre minus half a pixel. (\cref{ch:submission})
\item[Product parity] Checking that the product's code path reproduces the evaluator's scores on the same scenes. (\cref{ch:heavy})
\item[Progressive decoder] The FT09 decoder: SegFormer fusion at 1/4 resolution, then two x2 upsampling stages that each see a masked skip of the stem output. (\cref{ch:decoderwork})
\item[Project] A piece of work bound to exactly one Scene; it holds that work's Actions. (\cref{ch:product})
\item[Projection] The multiple of a direction vector closest to a given vector, found by minimising the squared distance; for a unit direction u it is (u.x) u. (\cref{ch:linalg})
\item[Prototype] A stored example in embedding space that an input is compared with; under level 1 by matching every library entry is one. (\cref{ch:namer})
\item[Pseudo-label] A training label produced by a model (the teacher) rather than a person. (\cref{ch:ftruns})
\item[PSF halo] The part of a small target's light that the sensor's point-spread function spreads into neighbouring pixels. (\cref{ch:uncertainty})
\item[QA PIXEL] Landsat's bit-packed quality layer; the builder rejects pixels with the fill, dilated cloud, cloud, cloud shadow or high-cirrus bits set. (\cref{ch:sensors})
\item[QR decomposition] Writing a matrix as an orthonormal-column matrix Q times a triangular R; Q spans the same space. (\cref{ch:decoderwork})
\item[Quality layer] A per-pixel mask the provider's processor computes with automatic tests, such as EnMAP's cloud or cirrus layer; not drawn by people. (\cref{ch:seghead})
\item[Quantile] The value below which a given fraction of a list lies; torch.quantile interpolates linearly between neighbouring sorted values. (\cref{ch:decoderwork})
\item[Query] The learned projection of a token that expresses what it is looking for; compared against every key. (\cref{ch:attention})
\item[Quotient rule] The derivative of a/b is (a'b - ab')/b squared. (\cref{ch:losses})
\item[R-theta branch] The unconstrained convolutional path in the HyperSL unmixing paper's reconstruction; it carried 29-56 percent of the reconstruction energy. (\cref{ch:hslunmix})
\item[Radiance] Power of light leaving a surface in one direction, per square metre, per steradian and per nanometre; what a sensor measures. (\cref{ch:light})
\item[Ramp] The schedule lambda(t) = lambda-max times min(1, t/R): a linear rise over R epochs, held constant within each epoch, then flat. (\cref{ch:losses})
\item[Ramped weight] A loss weight that grows over the first epochs (here the SAM weight, from 0 to 0.5 over 10 epochs); fine for training, wrong for ranking epochs. (\cref{ch:pretrain})
\item[Random inference mask] The alternative strategy: a random half of the tokens and its exact complement, which scatters token-boundary artefacts differently in every tile. (\cref{ch:inference})
\item[Rank] The number of independent directions a matrix's outputs fill; for a covariance, the number of non-zero eigenvalues. (\cref{ch:localrx})
\item[Rank (of a matrix)] The number of linearly independent columns of a matrix; the dimension of the set of vectors it can produce. (\cref{ch:compressor})
\item[Rank (of a tensor)] The number of axes a tensor has. (\cref{ch:shapes})
\item[Rank fusion] Combining several rank maps pixel by pixel by their mean (needs agreement) or their maximum (lets either scorer's false alarms through). (\cref{ch:scoringgains})
\item[Rank map] A score map converted to the fraction of valid pixels scoring at or below each pixel, so scorers with different scales can be combined. (\cref{ch:scoringgains})
\item[Rank of a matrix] The dimension of the set of outputs a linear map can produce; a map through 32 numbers has rank at most 32. (\cref{ch:decoderwork})
\item[Rank-nullity theorem] For a matrix with N columns, rank plus the dimension of the null space equals N. (\cref{ch:ceiling})
\item[read\_bands] The inferencer helper that keeps only bands within 15 nm of the trained-wavelength record and tells the model those wavelengths. (\cref{ch:shipping})
\item[Reading] How a model reads a scene's bands: its own native bands by wavelength, or a resample onto a fixed grid (product165, se202). (\cref{ch:runeval})
\item[Reading path] The chain of code that turns a stored scene into the cube a model reads; two paths for the same scene must give the same cube. (\cref{ch:sortbug})
\item[Realised contrast] The contrast z measured on the object after implanting, including the spot's own natural difference from its ring. (\cref{ch:benchmarks})
\item[Reaper] The routine that marks running jobs failed when their heartbeat is older than a timeout. (\cref{ch:product})
\item[Recall] For a class, the share of the pixels labelled with it that the classifier also calls it: diagonal entry over row total. (\cref{ch:hslcloud})
\item[Recall at a fixed false-alarm rate] The share of man-made plantings above the cut that calls a given share of natural plantings man-made; the namer's level-1 pass test uses ACE's rate. (\cref{ch:namer})
\item[Receptive field] The region of input pixels that can affect one output value. (\cref{ch:shapes})
\item[Reconstruction] The model's output, written x hat; its estimate of the recorded cube. (\cref{ch:recon})
\item[Reconstruction model] A function from an input cube to an output cube of the same shape, trained so that its output, the reconstruction, is close to its input on ordinary data. (\cref{ch:recon})
\item[Red edge] The sharp rise in vegetation reflectance between about 680 and 760 nm, caused by chlorophyll absorbing red and leaf structure scattering near-infrared light. (\cref{ch:light})
\item[Reduced stem] The reduced-attention stem: attention over fixed 40 nm wavelength bins plus three per-band paths; 14,721 compressor and 4,257 decoder parameters. (\cref{ch:specdecoder})
\item[Reduction LayerNorm] The LayerNorm applied to the summary tokens before they are projected into keys and values. (\cref{ch:esa})
\item[Reduction ratio] The factor R by which ESA shrinks the token grid in each direction before forming keys and values; 8, 4, 2 and 1 at the four stages. (\cref{ch:esa})
\item[Refine convolutions] The decoder's two 3 by 3 convolutions at quarter resolution, its only spatial layers; the second outputs 512 channels and has no activation. (\cref{ch:decoder})
\item[Reflectance] The fraction of arriving light that a surface returns at a given wavelength; unitless, usually between 0.02 and 0.30 for real ground. (\cref{ch:light})
\item[Regular expression] A text pattern, such as asphalt|concrete, matched against names; the family library's curation assigns families with them, first match wins. (\cref{ch:material})
\item[Release candidate] The model proposed for shipping; here FT04 with cloud exclusion on its own masks. (\cref{ch:ftruns})
\item[ReLU] Rectified Linear Unit, max(0, x). (\cref{ch:block})
\item[Remove tokens] Deleting hidden tokens from the sequence by a gather, so the encoder's blocks never see them. (\cref{ch:masking})
\item[Renormalisation] The rescaling of surviving softmax weights so they again sum to 1 after some positions are masked; ratios between survivors are unchanged. (\cref{ch:crossstem})
\item[Resample matrix] The (G, C) matrix that maps a spectrum on C scene bands to G grid wavelengths; its b-th column is the interpolation of the b-th unit vector. (\cref{ch:sortbug})
\item[Resampled band list] Values interpolated onto a band list other than the sensor's own, such as the 165-band grid. (\cref{ch:pretrain})
\item[Resampling] Estimating a spectrum's values at a fixed list of wavelengths from its values at the sensor's own wavelengths. (\cref{ch:sensors})
\item[Reshape] An operation that changes a tensor's shape while keeping its numbers in the same memory order. (\cref{ch:shapes})
\item[Residual] The recorded value minus the reconstructed value, x minus x hat, one number per band per pixel. (\cref{ch:recon})
\item[Residual (error)] The difference between the reconstruction and the truth, band by band, at one pixel. (\cref{ch:losses})
\item[Residual connection] Adding a sublayer's output to its input, x + f(x), so that the derivative is 1 + f' and an untouched path always exists; also called a skip connection. (\cref{ch:block})
\item[Residual covariance] The covariance matrix of the residuals of a scene's valid pixels, estimated in the product from up to 200,000 pixels with a fixed seed. (\cref{ch:residualmaha})
\item[Residual Mahalanobis distance] The squared Mahalanobis distance of a pixel's residual from the scene's mean residual, using the residual covariance; Indradhanu's default score since 2026-10-15. (\cref{ch:residualmaha})
\item[Resolved gain] A difference whose 95 percent paired-bootstrap interval lies wholly above zero; the repository's rule since 2026-10-21. (\cref{ch:bootstrap})
\item[Resolver] The static table of what each model architecture can do, combined with its manifest to build the model picker. (\cref{ch:product})
\item[Responsibility] The probability that a pixel came from one mixture component, given the pixel, by Bayes' rule. (\cref{ch:scoringgains})
\item[Responsibility-weighted Mahalanobis distance] The sum over components of each responsibility times the pixel's squared Mahalanobis distance to that component. (\cref{ch:scoringgains})
\item[Restore tokens] Putting the encoded kept tokens back at their positions by a scatter into zeros; hidden positions stay zero. (\cref{ch:masking})
\item[Resume mode] Whether a checkpoint restores the full training state (resume) or only the weights with a fresh optimiser and schedule (finetune). (\cref{ch:trainloop})
\item[Ring] The pixels inside the outer window but outside the guard: 61 x 61 minus 21 x 21 = 3,280 pixels by default. (\cref{ch:residualmaha})
\item[Ring fallback] The rule that a pixel whose ring holds fewer than 2 d valid pixels (60 at d = 30) is scored against the global mean instead. (\cref{ch:residualmaha})
\item[Ring mean] The mean of the valid pixels at Chebyshev distance 3 to 7 from a pixel (15 x 15 box minus 5 x 5), 200 pixels in the interior. (\cref{ch:namer})
\item[Robust loss] A loss whose pull from any single observation is bounded, so outliers cannot dominate it; L1 is robust, squared error is not. (\cref{ch:losses})
\item[ROC curve] The path of (FPR, TPR) points as the threshold slides from the top score to the bottom score. (\cref{ch:measures})
\item[Rotation] Swapping one block of training data for the next while the learning-rate schedule and optimiser state continue unbroken. (\cref{ch:pretrain})
\item[Row-major order] Memory order in which the last index changes fastest; offset of (a, b, d) in shape (A, B, D) is a B D + b D + d. (\cref{ch:sensors})
\item[Running mean] BatchNorm's stored estimate of a channel's mean, updated in training mode as 0.9 x old + 0.1 x batch mean. (\cref{ch:finetune})
\item[Running statistics] BatchNorm's exponential moving averages of batch mean and variance, kept as buffers and used in evaluation mode. (\cref{ch:compressor})
\item[RX detector] The Reed-Xiaoli detector: scores each pixel by its Mahalanobis distance from the scene mean. (\cref{ch:problem})
\item[SAM dark fade] An opt-in patch that multiplies the spectral angle by clip(mean reflectance / t, 0, 1) so that dark pixels such as open water do not dominate; superseded by residual Mahalanobis. (\cref{ch:residualmaha})
\item[SAM weight (lambda)] The multiplier on the spectral-angle term; it both schedules and converts radians into the scale of reflectance. (\cref{ch:losses})
\item[Sample-weighted average] An average over batches in which each batch's mean is weighted by the number of samples that survived filtering. (\cref{ch:trainloop})
\item[Sampled ROC] The product's ROC from evenly spaced percentiles of the scores, at most 256 of them (compute\_roc). (\cref{ch:measures})
\item[Sampling distribution] The spread of values a statistic would take across many fresh datasets of the same kind. (\cref{ch:bootstrap})
\item[Saturation (softmax)] The state in which one score is so far above the others that it takes almost all the weight and the softmax's derivatives are almost zero. (\cref{ch:attention})
\item[Scale and shift (gamma, beta)] BatchNorm's two learned parameters per channel, applied after standardising. (\cref{ch:compressor})
\item[Scale floor] The per-scene lower limit on b (its 5th percentile by default) that stops very confident pixels from inflating the score. (\cref{ch:uncertainty})
\item[Scale invariance] The property that multiplying a spectrum by a positive constant leaves its spectral angle to any other spectrum unchanged. (\cref{ch:losses})
\item[Scale invariance (of residual Mahalanobis)] The property that shrinking every residual by the same factor leaves every score unchanged, because the fitted covariance shrinks by the square of that factor. (\cref{ch:ftruns})
\item[Scale parameter] A positive number that stretches a distribution horizontally; for the Laplace it equals the mean absolute error. (\cref{ch:uncertainty})
\item[Scaled dot-product attention] softmax of Q times K-transpose divided by the square root of the head dimension, multiplied by V. (\cref{ch:attention})
\item[Scaled RM] The scorer scaled\_rm: residual Mahalanobis on the residual divided by the head's floored b (or its square root). (\cref{ch:uncertainty})
\item[Scatter] Writing values into named positions of a starting array; unnamed positions keep their starting value. (\cref{ch:masking})
\item[Scene] One complete satellite acquisition: one patch of ground at one moment, as one file or folder. (\cref{ch:problem})
\item[Scene (product)] One uploaded raw satellite acquisition, onboarded into a vendable. (\cref{ch:product})
\item[Scene adaptation] Briefly retraining a pretrained model on the scene being scored, then restoring its weights. (\cref{ch:scoringgains})
\item[Scene fingerprint] An identity of a scene's source files (file size and time, or SHA-256s) included in the cache key since commit 229d423. (\cref{ch:runeval})
\item[Scene-level uncertainty] Luck in which scenes are on the bench; not covered by the object bootstrap, so its intervals are lower bounds. (\cref{ch:bootstrap})
\item[Scene-major loop] Loading one scene at a time and running every model on it before the next scene, so memory holds one scene and one reconstruction. (\cref{ch:runeval})
\item[Scope] The set of surfaces and objects the product is meant to examine; excluding a class removes it from scope. (\cref{ch:runeval})
\item[Score map] The anomaly scores laid out on the pixel grid, one number per pixel. (\cref{ch:problem})
\item[Score tensor] The tensor of query-key scores, shape (B, h, N, N / R-squared); its size sets the memory cost of attention. (\cref{ch:esa})
\item[Score-map cache] Saved score maps under a key built from the scene, the model file and the scorer settings, so new measures need no model run. (\cref{ch:runeval})
\item[Score-map correction] Any post-processing applied to a finished anomaly score map, such as ranking each pixel within its edge-strength bin; none beat plain residual Mahalanobis. (\cref{ch:ftruns})
\item[Scorer] A rule turning a cube and its reconstruction into a score map, such as residual Mahalanobis or exclude. (\cref{ch:runeval})
\item[SDPA row chunking] Splitting the pixel rows of a fused-attention call into chunks of at most 32,768 to stay under CUDA's 65,535 launch-grid limit; exact because pixels are independent. (\cref{ch:crossstem})
\item[See-saw] FT02's per-sensor validation swinging between epochs, one sensor improving while others worsen, consistent with BatchNorm drift under sensor-pure batches. (\cref{ch:finetune})
\item[SegFormer] A hierarchical vision transformer with four scales and a light fusion decoder; the architecture behind Chakshu and Indradhanu. (\cref{ch:recon})
\item[SegFormer block] The repeating unit of the encoder: pre-norm efficient self-attention with a residual, then pre-norm Mix-FFN with a residual; shape (B, N, C) in and out. (\cref{ch:block})
\item[SegFormer-B0] The smallest published SegFormer configuration, with stage widths 32, 64, 160 and 256, which this encoder follows. (\cref{ch:encoder})
\item[Segmentation] Giving every pixel a label (water, cloud, ...) rather than one label per image. (\cref{ch:seghead})
\item[Segmentation head] A small network on the shared encoder's features that outputs one logit per class per pixel; in Indradhanu a SegFormerDecoder of width 128 with 5 outputs, 371,536 parameters. (\cref{ch:seghead})
\item[Segmentation weight] The factor seg\_weight multiplying the segmentation loss in joint training; 0.1 in FT03. (\cref{ch:finetune})
\item[segment\_mixed] The segmentation counterpart of forward\_mixed: one unmasked pass over the whole mixed batch, so BatchNorm sees every sensor. (\cref{ch:finetune})
\item[Selection effect] A shift in a population caused by which members survive a filter, not by any change to the members themselves. (\cref{ch:sortbug})
\item[Selection matrix] A matrix with one 1 per row that performs a gather by multiplication; its transpose performs the matching scatter into zeros. (\cref{ch:masking})
\item[Selection risk] The optimism that enters when a variant is chosen after seeing results on the same data that then grades it. (\cref{ch:uncertainty})
\item[Self-attention] Attention in which queries, keys and values all come from the same set of tokens. (\cref{ch:attention})
\item[Self-weight] How much a pixel contributes to its own reconstruction; a self-weight w shrinks every residual by the factor 1 minus w. (\cref{ch:recon})
\item[Sensor share] The fixed probability (or sub-batch fraction) with which each sensor is drawn: EnMAP 0.40, PRISMA 0.35, AVIRIS-NG 0.25. (\cref{ch:finetune})
\item[Sequence layout] The arrangement (B, N, C): a list of N tokens of C numbers each, used by attention, LayerNorm and linear layers. (\cref{ch:shapes})
\item[set\_wavelengths] The call that tells a spectral-stem model which wavelengths its next input holds; forgetting it leaves the previous scene's wavelengths in place. (\cref{ch:shipping})
\item[Shape] The list of sizes of an array's axes, such as (C, H, W). (\cref{ch:light})
\item[Shape-preserving model] A model whose output has exactly the input's shape, so that the two can be subtracted to give a residual. (\cref{ch:decoder})
\item[Shard] A tar archive packing many patches as named members, streamed sequentially; the webdataset convention. (\cref{ch:trainloop})
\item[Shard rotation] Training on a block of shards for a few epochs, then moving to a fresh block (Heavy: 200 shards every 5 epochs). (\cref{ch:heavy})
\item[Shift difference] Estimating the noise covariance as half the covariance of differences between neighbouring pixels, assuming neighbours share the same signal. (\cref{ch:localrx})
\item[Shift invariance] The property of the wavelength encoding that the similarity of two tags depends only on the difference of their wavelengths. (\cref{ch:ceiling})
\item[Shortwave infrared (SWIR)] Wavelengths from about 1000 to 2500 nm, where water, clays and minerals have absorption features. (\cref{ch:light})
\item[Shrink] One minus the ring's mean PSF weight (0.676 for a one-pixel object at sigma 1.4): the share of an implant's signal that survives as contrast against the ring. (\cref{ch:benchmarks})
\item[Shrinkage] Blending a covariance with its average eigenvalue times the identity, pulling all eigenvalues towards their mean; not used in this repository. (\cref{ch:rx})
\item[Siamese network] A network that runs two inputs through the same weights and compares the outputs; the namer's stem is shared this way by the pixel, its ring mean and every library spectrum. (\cref{ch:namer})
\item[Sigmoid] The function 1 / (1 + exp(-z)), which maps any number into (0, 1) and never reaches either end. (\cref{ch:specdecoder})
\item[Sign function] The function returning +1 for a positive number, -1 for a negative one and 0 for zero; the derivative of the absolute value away from zero. (\cref{ch:losses})
\item[Sign test] A test that counts in how many paired comparisons one method wins; with four scenes the smallest possible two-sided p-value is 0.125. (\cref{ch:residualmaha})
\item[Signal-to-noise ratio] Along a direction, the signal variance divided by the noise variance. (\cref{ch:localrx})
\item[Signal-to-noise ratio (SNR)] The filtered signal divided by the spread of the filtered noise, w-transpose t over the square root of w-transpose Sigma w. (\cref{ch:scoringgains})
\item[Signature implant] An additive implant x + a d, with d a real object's difference from its ring. (\cref{ch:benchmarks})
\item[Singular matrix] A square matrix with determinant zero, which has no inverse; a covariance is singular when the data do not vary at all along some direction. (\cref{ch:rx})
\item[Singular value decomposition] The factorisation X = U S V-transpose; the columns of V are the eigenvectors of X-transpose X, with eigenvalues the squared singular values. (\cref{ch:ftruns})
\item[Size score] Any score that measures how large a residual is (L1, MSE, SAM, combined), regardless of its direction. (\cref{ch:residualmaha})
\item[SKIP LOCKED] A Postgres row-locking clause that lets several workers claim different queued jobs without blocking or double-claiming. (\cref{ch:product})
\item[Skip path] The reduced stem's per-band route, the mean over valid bands of each value times W(lambda), straight to the 32 channels: a 1x1 compressor whose columns are computed from wavelength. (\cref{ch:specdecoder})
\item[Sliding window] Covering a large image with fixed-size tiles placed at regular offsets, here 128 x 128 tiles every 32 pixels. (\cref{ch:inference})
\item[Slope] Rise divided by run of a straight line: (y1 - y0)/(x1 - x0). (\cref{ch:bandpipeline})
\item[Snapped tile] The last tile along an axis, moved back so that it ends exactly at the scene edge; it overlaps its neighbour by more than the usual overlap. (\cref{ch:inference})
\item[Soft Dice loss] One minus a differentiable Dice coefficient computed from probabilities, with +1 smoothing; it gives each class a full-sized term however rare. (\cref{ch:seghead})
\item[Softmax] The function that turns a list of scores into positive weights summing to one, by exponentiating each score and dividing by the total. (\cref{ch:attention})
\item[Span] All weighted sums of a set of vectors. (\cref{ch:decoderwork})
\item[Spatial completion] Predicting a pixel's spectrum from other pixels, its neighbours; how Indradhanu reconstructs a hidden pixel. (\cref{ch:hsldetect})
\item[Spatial layout] The arrangement (B, C, H, W): a picture with C channels, used by convolutions. (\cref{ch:shapes})
\item[Spatial matched filter] A blur shaped like a target's point-spread footprint, applied to the residual before whitening (scoring-gains idea A1). (\cref{ch:scoringgains})
\item[Spatial-reduction attention] The general name, from the Pyramid Vision Transformer, for attention whose keys and values are computed on a downsampled grid. (\cref{ch:esa})
\item[Spearman rank correlation] The Pearson correlation of two lists' ranks; 1 for the same order, -1 for the reverse, near 0 for unrelated orders. (\cref{ch:uncertainty})
\item[Spectral angle] The angle between two spectra treated as vectors, arccos of their dot product over the product of their lengths, in radians; it measures shape and ignores brightness. (\cref{ch:losses})
\item[Spectral angle mapper] Matching by the angle between a pixel's spectrum and a library spectrum, which ignores overall brightness. (\cref{ch:material})
\item[Spectral Angle Mapper (SAM)] The remote-sensing method that classifies a spectrum by its spectral angle to library spectra; in this book also the name of the angle used as a loss. (\cref{ch:losses})
\item[Spectral compressor] Indradhanu's first learned module: a 1x1 convolution from 165 bands to D = 32 channels followed by BatchNorm, applied at every pixel. (\cref{ch:compressor})
\item[Spectral decompressor] The 1x1 convolution from 32 channels back to 165 bands at the end of Indradhanu, with no normalisation and no activation. (\cref{ch:compressor})
\item[Spectral family] In the code, the detector a band came from (the SpectralFamily enum: VNIR, SWIR, PANCHROMATIC); not the five grid segments. (\cref{ch:sensors})
\item[Spectral floor] The part of a spectrum outside the span of the spectral output layer's curves; no decoder can remove it. (\cref{ch:decoderwork})
\item[Spectral library] A collection of laboratory reflectance spectra of known materials; the product uses a curated subset of USGS splib07. (\cref{ch:material})
\item[Spectral namer] The two-tower network built on 2026-10-31 (app/spectral\_match/namer/) that names flagged pixels: P(man-made), a man-made class and a fill estimate. (\cref{ch:namer})
\item[Spectral response function (SRF)] The curve describing how strongly a band responds at each wavelength; a band's value is a weighted average of the spectrum with these weights. (\cref{ch:light})
\item[Spectral self-completion] Training a pixel to fill its own hidden bands from its other bands; rejected because an anomaly could then rebuild itself. (\cref{ch:specdecoder})
\item[Spectral theorem] Every real symmetric matrix can be written U Lambda UT with orthogonal unit eigenvectors in U and real eigenvalues in Lambda. (\cref{ch:linalg})
\item[Spectral tokenizer] HyperSL's front end: three 1-D convolutions over band index (kernels 5, 9, 13, stride 1) fused by a 1 x 1 convolution, giving one 256-number token per band. (\cref{ch:hslwhat})
\item[Spectral-stem model] The wavelength-aware variant of Indradhanu (V2) that replaces the fixed 165-band compressor and decompressor; same encoder and decoder classes. (\cref{ch:params})
\item[SpectralEarth] DLR's collection of 539,755 EnMAP patches, 128 by 128 pixels on 202 bands, published for pretraining hyperspectral models; Indradhanu's pretraining set. (\cref{ch:pretrain})
\item[Spectrum] One pixel's values across all bands, read in wavelength order; treated as a vector in R to the C. (\cref{ch:light})
\item[Split by site] Dividing data into training and test so that every image of one place lands on the same side, preventing the test score from flattering the model. (\cref{ch:pretrain})
\item[Square-root stretch] Displaying the square root of a score so that a heavy-tailed distribution of squared distances spreads across the colour scale. (\cref{ch:inference})
\item[Squeeze-and-excitation] A per-image (here per-patch) channel gate computed from a pooled summary; it gives every pixel of the patch the same recipe. (\cref{ch:specdecoder})
\item[Stable softmax] Computing the softmax after subtracting the largest score from all scores, which leaves the result unchanged and prevents overflow. (\cref{ch:attention})
\item[Stage] One level of the encoder: a patch embedding, a number of transformer blocks and a LayerNorm, at one grid size and width. (\cref{ch:encoder})
\item[Stage 3] The reference evaluation of 2026-10-17 (configs/eval/stage3.json) that reproduced earlier numbers and exposed bench-v1's bias. (\cref{ch:runeval})
\item[Standard deviation] The square root of the average squared distance from the mean; the typical spread of a set of numbers, in their own units. (\cref{ch:problem})
\item[Standard error] The standard deviation of a statistic's sampling distribution; sigma / sqrt(n) for the mean of n independent values. (\cref{ch:bootstrap})
\item[Standard normal] The Gaussian with mean 0 and standard deviation 1. (\cref{ch:rx})
\item[Statistical ensembler] A notebook-only detector that fuses global and local RX score maps by their ranks. (\cref{ch:localrx})
\item[Step (iteration)] One update of the parameters by the optimiser. (\cref{ch:trainloop})
\item[Stochastic gradient descent] Gradient descent using batch gradients, which are cheap, noisy estimates of the full gradient. (\cref{ch:trainloop})
\item[Strength] The single number (a, f or s) that scales an implant; solved so the realised contrast hits its target. (\cref{ch:benchmarks})
\item[Strict loading] The default behaviour of PyTorch's load\_state\_dict: every saved tensor must find a slot of the same name and shape, or loading fails with an error. (\cref{ch:inference})
\item[Stride] In local RX, scoring only every n-th row and column and interpolating the rest. (\cref{ch:localrx})
\item[Stride (of a convolution)] How far the kernel jumps between neighbouring output positions. (\cref{ch:shapes})
\item[Stride (of a tensor)] The memory distance between neighbouring elements along one axis. (\cref{ch:shapes})
\item[Stride (sliding window)] The distance between the starts of neighbouring tiles; 32 pixels for Indradhanu at inference. (\cref{ch:inference})
\item[Strided convolution] A convolution whose window moves more than one position at a time; with kernel equal to stride the windows tile the grid without overlap. (\cref{ch:esa})
\item[Sub-pixel convolution] A convolution that outputs r squared channels per output channel so that a pixel shuffle can turn them into an r by r block of pixels. (\cref{ch:decoder})
\item[Sublayer] One of the two halves of a block, attention or Mix-FFN, each with its own LayerNorm, dropout and residual. (\cref{ch:block})
\item[Summary token] One of the N / R-squared reduced tokens of ESA, each standing for an R-by-R block of the token grid; there are 16 at every stage for a 128-by-128 patch. (\cref{ch:esa})
\item[Summed-area table] Another name for an integral image. (\cref{ch:residualmaha})
\item[Surface-driven label] A cirrus label whose pixels differ from clear land mainly in the shortwave infrared, so it follows the ground rather than a cloud. (\cref{ch:seghead})
\item[Symmetric matrix] A square matrix equal to its own transpose, such as every covariance matrix. (\cref{ch:linalg})
\item[Target / ring ratio] Median b on a target's pixels divided by median b on pixels 4 to 8 city-block steps away; above 1 means the head expects more error on the target than on its own surroundings. (\cref{ch:uncertainty})
\item[Teacher-student training] Training a student model to match the output of a teacher that sees more; here the teacher would be an EMA of the student. (\cref{ch:ftruns})
\item[Template (file)] A dictionary, injected into a FileHelper, mapping each logical component to where it lives in the file. (\cref{ch:sensors})
\item[Tensor] In this book, a multi-index array of numbers. (\cref{ch:light})
\item[The 39] The 39 of 909 real labelled objects lying wholly inside FT03's water / cloud mask; 35 are on water or shorelines, 4 under haze in set4. (\cref{ch:runeval})
\item[Thermal sensor] A sensor recording emitted heat in one thermal-infrared band, giving a cube with C = 1. (\cref{ch:light})
\item[Thick-cloud rule] Since 2026-10-29, V2 excludes only the head's thick-cloud output, not thin cloud, because the Heavy model's thin-cloud output fires on haze. (\cref{ch:shipping})
\item[Thin-cloud channel] The segmentation head's third output; it fires on haze, which made it delete targets under set4's haze. (\cref{ch:heavy})
\item[Threshold] A score above which a pixel is flagged; chosen from the false-alarm rate one accepts. (\cref{ch:rx})
\item[Tier] Subtle, medium or easy: contrast z below 0.502, up to 0.934, or above; the terciles of the 909 real objects. (\cref{ch:benchmarks})
\item[Tile validity filter] The rule that skips any tile with less than 10 percent valid pixels at inference (training used 40 percent). (\cref{ch:inference})
\item[Tiled segmentation] Running the segmentation head in overlapping 128 px windows at stride 64 and averaging the probabilities where windows overlap. (\cref{ch:shipping})
\item[Tirohita] Antardhana's L1 sibling, smaller per stage, deeper, and hiding 50 to 75 percent; 329,665 parameters. (\cref{ch:convae})
\item[To-channel layer] The Linear(32, 1) shared by all 32 answers that turns each answer vector into one output channel. (\cref{ch:crossstem})
\item[Token] A small block of the image, 4 by 4 pixels at Indradhanu's first stage, summarised as one vector. (\cref{ch:recon})
\item[Token bottleneck] Stage 1's squeeze of 4 x 4 pixels x 32 channels = 512 numbers into one 32-wide token vector. (\cref{ch:decoderwork})
\item[Token grid] The H' by W' arrangement of tokens produced by a patch embedding; 32 by 32 at stage 1 for a 128 by 128 patch. (\cref{ch:patchembed})
\item[Token mask] The (B, N) array of token validities: whether each token holds real data. A fact about the sensor, not a choice. (\cref{ch:masking})
\item[Token straddling] An object that crosses stage-1 token borders, so in each masked pass part of it stays visible and helps reconstruct the hidden part, lowering its residual. (\cref{ch:chakshu})
\item[Token validity] A token is valid when its valid fraction is strictly greater than one half; exactly half counts as invalid. (\cref{ch:masking})
\item[Token width] The number of components d in each token; 32 in the spectral stem. (\cref{ch:ceiling})
\item[Top-of-atmosphere reflectance] pi L / (E0 cos of the solar zenith): a crude conversion of radiance that removes the sun but not the atmosphere. (\cref{ch:benchmarks})
\item[Topological order] An ordering of a DAG's nodes in which every node comes after all of its parents; found with Kahn's algorithm. (\cref{ch:product})
\item[torchinfo] A PyTorch library that prints a per-layer table of shapes, parameter counts and multiply-adds. (\cref{ch:params})
\item[Total variation] The sum over side-by-side pixel pairs of the absolute difference of their values; as a penalty it shrinks every target by an amount proportional to perimeter over area. (\cref{ch:hslunmix})
\item[Tower] One encoder in a model that has several side by side; the namer has a pixel tower and a library tower. (\cref{ch:namer})
\item[TPR at a fixed FPR] The best detection rate at any exact ROC point whose false-alarm rate is at most the given value. (\cref{ch:measures})
\item[Trace] The sum of a square matrix's diagonal entries; equal to the sum of its eigenvalues. (\cref{ch:rx})
\item[Trained-band mask] The rule that keeps only bands within 15 nm of a wavelength the model has trained on, guarding against extrapolation. (\cref{ch:specdecoder})
\item[Trained-wavelength record] The buffer trained\_wl listing every wavelength the model has trained on; it grows during fine-tuning and decides which bands are accepted at load time. (\cref{ch:finetune})
\item[Training mode] The mode in which BatchNorm normalises with the current batch's statistics and updates its running statistics. (\cref{ch:compressor})
\item[Transpose] A matrix with rows and columns swapped; the transpose of a product reverses the order, (AB)T = BT AT. (\cref{ch:linalg})
\item[Transpose (axes)] Reordering an array's axes without changing any value, such as BIL to BSQ. (\cref{ch:sensors})
\item[Transposed convolution] The learnable upsampling layer that spreads each input value over a kernel-sized stretch of the output; with kernel 4, stride 2, padding 1 it doubles height and width. (\cref{ch:convae})
\item[Trapezoid rule] Area of a straight segment as width times the average of its two heights; summed over segments it gives the area under a piecewise-linear curve. (\cref{ch:measures})
\item[Trim fraction] The share tau of the largest per-pixel errors a trimmed loss discards; 0.05 for Chakshu, 0 for Indradhanu. (\cref{ch:chakshu})
\item[Trimmed loss] A loss that sorts the per-pixel errors and averages only the smallest fraction 1 minus tau of them, so the worst (presumed anomalous) pixels give no gradient; Chakshu trims 5 percent. (\cref{ch:chakshu})
\item[True-positive rate (TPR)] The share of targets flagged at a threshold; also called the detection rate. (\cref{ch:measures})
\item[Truncation (PCA)] Keeping only the d directions with the largest eigenvalues (d = 30 here) and discarding the rest, which stabilises the score at the cost of sensitivity in the quietest directions. (\cref{ch:residualmaha})
\item[Two-level namer] A naming design (2026-10-31; built the same day as the spectral namer, \cref{ch:namer}): level 1 asks whether nature explains a pixel, level 2 which man-made class it is; P(class) = P(man-made) x P(class | man-made). (\cref{ch:material})
\item[Two-pass inference] Running the model twice with complementary masks and keeping each pixel's prediction from the pass in which it was hidden. (\cref{ch:recon})
\item[Two-pass masking] Running the model twice on a tile with complementary token masks and taking each pixel's prediction from the pass in which it was hidden, so that no pixel is ever predicted while visible. (\cref{ch:inference})
\item[Two-pass validation] Validation that hides a random half of the tokens, then the complementary half, so every valid pixel is graded exactly once from a prediction that did not see it. (\cref{ch:losses})
\item[Two-point form] The equation of a line through two known points: y = y0 + slope times (x - x0). (\cref{ch:bandpipeline})
\item[Two-tower model] A model with two encoders whose outputs are compared, usually by a cosine; also called a dual encoder. (\cref{ch:namer})
\item[Unbiased variance] A variance that divides by M minus 1; PyTorch feeds it to BatchNorm's running variance. (\cref{ch:compressor})
\item[Uncertainty head] A small third decoder on Indradhanu's encoder features that predicts, per pixel, the log of the reconstruction error the model expects there (FT06, 95,632 parameters). (\cref{ch:uncertainty})
\item[Undefined behaviour (by contract)] Output that a function's documentation declares meaningless for some inputs; whatever a given version returns cannot be relied on. (\cref{ch:sortbug})
\item[Union bound] The chance that at least one of several events happens is at most the sum of their chances. (\cref{ch:material})
\item[Unit vector] A vector of length 1, used to represent a pure direction. (\cref{ch:linalg})
\item[Unpredictability] Error that remains because a hidden pixel's value is not fixed by its context; it is what makes targets detectable. (\cref{ch:decoderwork})
\item[Upscale factor] The factor r of a pixel shuffle; 4 here, taking 32 by 32 to 128 by 128. (\cref{ch:decoder})
\item[UUID] A 128-bit random identifier; on the wire always written with a type prefix such as scene\_ or action\_. (\cref{ch:product})
\item[Valid fraction] The proportion of a token's 16 pixels that are valid, computed as the average of the 0/1 pixel mask over the token's block. (\cref{ch:masking})
\item[Valid-band loss] A reconstruction loss computed only over bands that were measured; dead and padded bands are excluded from the L1 mean and zeroed before SAM. (\cref{ch:specdecoder})
\item[Validity cube] A 0/1 array of the cube's shape marking which band values of which pixels are real measurements. (\cref{ch:light})
\item[Validity mask] An array on the image grid holding 1 for a genuine measurement and 0 for dead detectors, cloud or off-swath pixels. (\cref{ch:problem})
\item[Value] The learned projection of a token that it hands over, in proportion to its attention weight, when another token attends to it. (\cref{ch:attention})
\item[Value stretch] The Linear(1, d) layer that turns one reflectance number into a d-number vector; 64 parameters at d = 32. (\cref{ch:ceiling})
\item[Vanishing gradient] The shrinking of the error signal to almost nothing after it is multiplied by many factors smaller than one on its way back through a deep network. (\cref{ch:block})
\item[Variance] The average squared distance of a random number from its mean; its square root is the standard deviation. (\cref{ch:attention})
\item[Variance of a difference] Var(A - B) = Var A + Var B - 2 Cov(A, B); small when A and B move together. (\cref{ch:bootstrap})
\item[VCA] Vertex component analysis: finds endmembers as the most extreme pixels of the data cloud. (\cref{ch:hslunmix})
\item[Vector] An ordered list of real numbers, written as a column; a spectrum on the 165-band grid is a vector in R to the 165. (\cref{ch:light})
\item[Vendable] The standard in-memory form of a scene: cube in physical units, validity cube and per-band metadata; consumed by everything downstream. (\cref{ch:product})
\item[View] A tensor operation that reinterprets the shape of stored numbers without copying them, used to split 32 numbers into 4 heads of 8. (\cref{ch:crossstem})
\item[Visible neighbours] Of the 8 tokens around a hidden token, those the encoder can see; expected 2.81 for random single tokens and 1.77 for 2 x 2 blocks at ratio 0.65. (\cref{ch:decoderwork})
\item[Vision Transformer (ViT)] A transformer for images that cuts the image into non-overlapping patches and embeds each by a linear projection. (\cref{ch:patchembed})
\item[VNIR and SWIR detectors] The two detector arrays of a hyperspectral instrument, for the visible/near-infrared and the shortwave-infrared ranges. (\cref{ch:sensors})
\item[Warm start] Initialising a component from previously trained weights rather than at random, as FT03's head from FT01's probe. (\cref{ch:finetune})
\item[Warmup] A schedule phase that raises the learning rate gradually from a small value at the start of training; described by warmup-epochs in this repository's config but never implemented. (\cref{ch:trainloop})
\item[Water guard] The rule that a confident water pseudo-label over a mean 800-900 nm reflectance above 0.15 becomes a known not-water. (\cref{ch:ftruns})
\item[Water vapour column] The mass of water in the air above each square centimetre of ground, in g/cm2; 0.75 g/cm2 would condense to a 7.5 mm layer. (\cref{ch:seghead})
\item[Wavelength] The distance between successive crests of a light wave, written lambda and measured in nanometres; different wavelengths are different colours, most of them invisible. (\cref{ch:light})
\item[Wavelength bin] One cell of the fixed grid of 40 nm bins from 400 nm; 53 bins up to 2500 nm, the same on every sensor. (\cref{ch:specdecoder})
\item[Wavelength encoding] A fixed vector of sines and cosines of 2 pi lambda / P over 16 geometric periods P from 10 to 10,000 nm; sines first, then cosines. (\cref{ch:ceiling})
\item[Wavelength query] A decoder query built from the wavelength encoding of a band to be reconstructed, through a Linear(32, 32). (\cref{ch:specdecoder})
\item[Wavelength statistics] WavelengthStats: mean and standard deviation curves over wavelength, measured once on SpectralEarth and read at any band by linear interpolation. (\cref{ch:specdecoder})
\item[Wavelength-linear decoder] WavelengthLinearDecoder: each requested band is U(lambda) . z + c(lambda), with U and c from one hypernetwork. (\cref{ch:specdecoder})
\item[Wavelength-reading model] A model with a spectral stem that reads a scene's own bands at their stated wavelengths (F001 and the FT runs), so it needs no resample. (\cref{ch:sortbug})
\item[WavelengthStats] The normalisation of Indradhanu V2: mean and spread stored as curves over wavelength and read at each sensor's own band wavelengths. (\cref{ch:bandpipeline})
\item[Weight] A learned multiplier applied to one input of a neuron or one entry of a kernel. (\cref{ch:shapes})
\item[Weight decay] Shrinking every weight towards zero by a fixed fraction each step, a form of regularisation. (\cref{ch:trainloop})
\item[Weight sharing] Using the same kernel weights at every position, which makes a convolution's parameter count independent of image size. (\cref{ch:shapes})
\item[Weighted mean] A sum of values times weights divided by the sum of the weights. (\cref{ch:ftruns})
\item[Whitened coordinate] One component of a whitened vector; the sum of their squares is the Mahalanobis distance. (\cref{ch:residualmaha})
\item[Whitening] Rotating vectors onto the principal directions and dividing each coordinate by its standard deviation, so the cloud has identity covariance. (\cref{ch:linalg})
\item[Width] The number of channels a layer carries. (\cref{ch:heavy})
\item[Width weighting] Weighting each band by the nanometres of spectrum it stands for, so an average over bands approximates an integral over wavelength. (\cref{ch:pretrain})
\item[Width-squared scaling] Because most weights sit in C x C matrices, multiplying width by k multiplies parameters by about k squared (8.86 for k = 3 at stage 1). (\cref{ch:heavy})
\item[Wire format] A data layout that stored files and saved models depend on by position, so that changing it silently invalidates them. (\cref{ch:sensors})
\item[With replacement] Drawing so that each pick can be any item, including ones already picked; some items appear several times, some not at all. (\cref{ch:bootstrap})
\item[Worker] The process that claims jobs from the jobs table and runs the analysis. (\cref{ch:product})
\item[z-score] A value expressed as standard deviations from its band's mean: z = (x - mean)/std. (\cref{ch:bandpipeline})
\item[Zero padding] Treating positions outside the grid as zeros so that a convolution keeps the grid size; it also makes border positions distinguishable. (\cref{ch:block})
\end{description}
